% concatenated from MDP3rdKind_plain.tex
\documentclass[reqno, 11pt]{amsart}%
\usepackage{amsaddr} 
\usepackage{amssymb}
\usepackage[nesting]{hyperref}
\usepackage[pdftex]{graphicx}
\usepackage{listings}
\usepackage{multirow}
\usepackage{placeins}
\usepackage{color}
%\usepackage{subfigure}
\usepackage{lscape}
\usepackage{dsfont}
\usepackage{amsmath}
\usepackage{amsfonts}
\usepackage{verbatim}
\usepackage{accents} 
\usepackage{todonotes}
\usepackage{mathtools}
\usepackage{bbm}
\usepackage{natbib}
\usepackage{lipsum}
%\usepackage{subcaption}
\usepackage{caption}
\usepackage{subcaption}
\usepackage{afterpage}

%\usepackage{blindtext}

% Any additional packages needed should be included after jmlr2e.
% Note that jmlr2e.sty includes epsfig, amssymb, natbib and graphicx,
% and defines many common macros, such as 'proof' and 'example'.
%
% It also sets the bibliographystyle to plainnat; for more information on
% natbib citation styles, see the natbib documentation, a copy of which
% is archived at http://www.jmlr.org/format/natbib.pdf

% Available options for package jmlr2e are:
%
%   - abbrvbib : use abbrvnat for the bibliography style
%   - nohyperref : do not load the hyperref package
%   - preprint : remove JMLR specific information from the template,
%         useful for example for posting to preprint servers.
%
% Example of using the package with custom options:
%
% \usepackage[abbrvbib, preprint]{jmlr2e}

%\usepackage{jmlr2e}

\usepackage[ruled,vlined]{algorithm2e}

\usepackage{xcolor}
\usepackage{upgreek}
\usepackage{float}

% BOLD UPPER CASE
\newcommand\bA{\mathbf A}
\newcommand\bB{\mathbf B}
\newcommand\bF{\mathbf F}
\newcommand\bN{\mathbf N}
\newcommand\bU{\mathbf U}
\newcommand\bW{\mathbf W}
\newcommand\bX{\mathbf X}
\newcommand\bY{\mathbf Y}
\newcommand\bZ{\mathbf Z}

% BOLD LOWER CASE
\newcommand\ba{\mathbf a}
\newcommand\br{\mathbf r}
\newcommand\bu{\mathbf u}
\newcommand\bv{\mathbf v}
\newcommand\bw{\mathbf w}
\newcommand\by{\mathbf y}
\newcommand\bz{\mathbf z}

% BOLD GREEK
\def\balpha{\boldsymbol{\alpha}}
\def\bbeta{\boldsymbol{\beta}}
\def\bepsilon{\boldsymbol{\epsilon}}
\def\bgamma{\boldsymbol{\gamma}}
\def\bmu{\boldsymbol{\mu}}
\def\bnu{\boldsymbol{\nu}}
\def\btheta{\boldsymbol{\theta}}
\def\bpi{\boldsymbol{\pi}}
\def\bvarphi{\boldsymbol{\varphi}}


% CALLIGRAPHIC
\newcommand\cA{\mathcal A}
\newcommand\cB{\mathcal B}
\newcommand\cC{\mathcal C}
\newcommand\cD{\mathcal D}
\newcommand\cE{\mathcal E}
\newcommand\cF{\mathcal F}
\newcommand\cG{\mathcal G}
\newcommand\cI{\mathcal I}
\newcommand\cK{\mathcal K}
\newcommand\cL{\mathcal L}
\newcommand\cM{\mathcal M}
\newcommand\cN{\mathcal N}
\newcommand\cP{\mathcal P}
\newcommand\cS{\mathcal S}
\newcommand\cT{\mathcal T}
\newcommand\cU{\mathcal U}
\newcommand\cV{\mathcal V}
\newcommand\cX{\mathcal X}
\newcommand\cY{\mathcal Y}


% MATHBB
\def\AA{\mathbb{A}}
%\newcommand\AA{\mathbb A}
\newcommand\EE{\mathbb E}
\newcommand\FF{\mathbb F}
\newcommand\JJ{\mathbb J}
\newcommand\PP{\mathbb P}
\newcommand\NN{\mathbb N}
\newcommand\RR{\mathbb R}
\newcommand\TT{\mathbb T}
\newcommand\VV{\mathbb V}
\newcommand\ZZ{\mathbb Z}

\def \kP{\overline{\textsf{\sf P}}}
\def \kS{\overline{S}}
\def \kA{\overline{A}}
\def \rs{\mathrm s}
\def \ra{\mathrm a}

%\newtheorem{thm}{Theorem}[section]
%\newtheorem{lem}[thm]{Lemma}%[section]
%\newtheorem{prop}[thm]{Proposition}%[section]
%\newtheorem{cor}[thm]{Corollary}
%%\theoremstyle{definition}
%\newtheorem{defn}[thm]{Definition}%[section]
%%\newtheorem{conj}{Conjecture}[section]
%\newtheorem{exmp}[thm]{Example}%[section]
%%\theoremstyle{remark}
%\newtheorem{remark}[thm]{Remark}%[section]
%\newtheorem{prob}[thm]{Problem}%[section]

\newtheorem{theorem}{Theorem}
\newtheorem{corollary}[theorem]{Corollary}
\newtheorem{lemma}[theorem]{Lemma}
\newtheorem{proposition}[theorem]{Proposition}
\theoremstyle{definition}
\newtheorem{example}[theorem]{Example}
\newtheorem{remark}[theorem]{Remark}
\newtheorem{definition}[theorem]{Definition}
\newtheorem{assumptions}[theorem]{Assumption}
\newtheorem{prob}[theorem]{Problem}


\newtheorem{assumption}{Assumption}  % independent global counter
\setcounter{assumption}{0}          

% start from 0%\newtheorem{note}{Note}
\textheight=24cm \textwidth = 16cm \topmargin= -1cm \oddsidemargin
0mm \evensidemargin 0mm
\newcommand{\BlackBoxes}{\global\overfullrule5pt}
\newcommand{\NoBlackBoxes}{\global\overfullrule0pt}
\BlackBoxes
\setlength{\overfullrule}{0pt}

% Heading arguments are {volume}{year}{pages}{date submitted}{date published}{paper id}{author-full-names}

\usepackage{lastpage}
%\jmlrheading{xx}{xx}{1-\pageref{LastPage}}{1/26; Revised xx}{xx}{21-0000}{B\"auerle and Vasileiadis}

% Short headings should be running head and authors last names

%\ShortHeadings{Markov Decision Processes of the Third Kind}{B\"auerle and Vasileiadis}
%\firstpageno{1}

\begin{document}

\title[MDP of the Third Kind]{Markov Decision Processes of the Third Kind: Learning Distributions by Policy Gradient Descent}

\author[N. \smash{B\"auerle}]{Nicole B\"auerle}
\address[N. B\"auerle]{Department of Mathematics,
Karlsruhe Institute of Technology (KIT), D-76128 Karlsruhe, Germany, \href{mailto:nicole.baeuerle@kit.edu}{nicole.baeuerle@kit.edu}}

\author[A. \smash{Vasileiadis}]{Athanasios Vasileiadis$^*$}
\address[A. Vasileiadis]{Department of Mathematics,
Karlsruhe Institute of Technology (KIT), D-76128 Karlsruhe, Germany, \href{mailto:nicole.baeuerle@kit.edu}{athanasios.vasileiadis@kit.edu}, \footnote{$^*$Corresponding author}}



%\author{\name Author One \email one@stat.washington.edu \\
%       \addr Department of Statistics\\
%       University of Washington\\
%       Seattle, WA 98195-4322, USA
%       \AND
%       \name Author Two \email two@cs.berkeley.edu \\
%       \addr Division of Computer Science\\
%       University of California\\
%       Berkeley, CA 94720-1776, USA}

%\editor{My editor}



\begin{abstract}%   <- trailing '%' for backward compatibility of .sty file
The goal of this paper is to analyze distributional Markov Decision Processes as a class of control problems in which the objective is to learn policies that steer the distribution of a cumulative reward toward a prescribed target law, rather than optimizing an expected value or a risk functional. To solve the resulting distributional control problem in a model-free setting, we propose a policy-gradient algorithm based on neural-network parameterizations of randomized Markov policies, defined on an augmented state space and a sample-based evaluation of the characteristic-function loss. Under mild regularity and growth assumptions, we prove convergence of the algorithm to stationary points using stochastic approximation techniques. Several numerical experiments illustrate the ability of the method to match complex target distributions, recover classical optimal policies when they exist, and reveal intrinsic non-uniqueness phenomena specific to distributional control.
\end{abstract}

\maketitle

\makeatletter \providecommand\@dotsep{5} \makeatother

\vspace{0.5cm}
\begin{minipage}{14cm}
{\small
\begin{description}
\item[\rm \textsc{ Key words}]
{\small Distributional MDPs, Policy Gradient, Randomized controls }
%\item[\rm \textsc{JEL subject classifications}] {...}
%{\small 91A35, 91A16, 91G20}

\end{description}
}
\end{minipage}\\

\section{Introduction}
Since the introduction of Reinforcement Learning (RL), tremendous progress has been achieved in designing algorithms that allow agents to learn optimal behavior through interaction with uncertain environments. Markov Decision Processes (MDPs) has been so far the theoretical framework to model these problems, \cite{Puterman1994MDP, Bertsekas-Shreve,Powell2011ADP,BauerleRieder}. When all the involved functions i.e. the dynamics and rewards/costs, are known and the spaces finite, we can use exact dynamic programming methods usually in a tabular form. When the spaces are continuous we need some form of function approximation. For what concerns this present article, RL is what we do when we only observe samples from the states actions and rewards, without access to any part of the theoretical model.     
 
 Classical RL methods focus on estimating the expected cumulative return, the classical value function of control, which has led to impressive successes across domains ranging from notoriously hard games like chess and GO to modern life robotics. We call this mathematical formulation \emph{MDPs of the first kind}. However, not all problems we are facing can be solved by the classical value functions, for example portfolio selection under model uncertainty, robust control and in general risk aware optimization, \cite{ContModelUncertainty,bauerleglaunerDistributionallyRobustMDPs,PrashanthFu2022RiskSensitiveRL,CoacheJaimungalRLconvexriskmeasures,coache2025robustreinforcementlearningdynamic,BauerleJaskiewiczMDPRiskSensitiveOverview,Guin_2026} call for a different formulation under coherent risk measures and a time consistent dynamic programming principle, \cite{ArtznerCoherentMeasuresofRist}. We call collectively this mathematical formulation \emph{MDPs of the second kind}. Finally, motivated by recent advances in finance, \cite{LassanceVrins2023TargetDistribution, ColaneriEisenbergSalterini2023TerminalDist}  and AI (especially generative AI), we are interested in optimal controls that stir the whole distribution of (discounted) cumulative rewards towards a certain target distribution. %Following \cite{bauerle2025distributionalbellmanequation}
 We call these (Distributional) MDPs,  \emph{MDPs of the third kind}. In this article we restrict to problems with a finite time hoizon.
 
 In complete analogy, Distributional Reinforcement Learning (DRL) addresses the problem in a practical way, using only samples by modeling the entire distribution of returns, providing a richer representation that can enhance stability, performance, and risk-sensitive decision making. 

\subsection{Motivation}
\label{se:motivation}
To motivate our results we start by giving an answer to a simple variant of our problem that we can actually solve analytically, so we get a flavor of what to expect in more complicated and general cases that we treat in the rest of the article. 

\begin{prob}
\label{prob:informal}
	Find optimal controls that shift the distribution of a terminal cumulative (discounted) reward (cost) such that it matches a predefined target. 
\end{prob}

On a compact space, the Fourier representation of a probability measure is a sequence of coefficient and matching this is a finite dimensional problem. In particular we can match directly by identifying the coefficients of sequence.  

To illustrate this, consider $\mathbb{T}^{1} = \mathbb{R}/\mathbb{Z}$, the one dimensional torus endowed with its Borel $\sigma$-algebra.
For a finite horizon $T \in \mathbb{N}$ consider the Markov Decision Process
\begin{equation*}
%\label{eq:def_lin_dynamics}
s_{t+1} = \bigl(s_{t} + a_{t} + \varepsilon_{t+1}\bigr) \bmod 2\pi,
\qquad t = 0,\dots,T-1,
\end{equation*}
where $s_0 \in \mathbb{T}^{1}$ is an $\mathcal F_0$-measurable initial condition with distribution $\mu_0$. The reward is given by $r_t(s_t,a_t)=0$ for $t=0,1,\ldots,T-1$ and $r_T(s_T,a_T)=s_T.$  Next,  $a_{t}\in\mathcal A\subset\mathbb{T}^{1}$ is the action chosen by the decision maker at stage $t$ which may depend on the observed history of the process. We consider randomized policies here.
%also the case of stochastic controls sampled from a distribution $\nu$ (possibly also with a density). 
Finally, $(\varepsilon_{t})_{t\ge 1}$ is an i.i.d. white noise sequence on $\RR$, independent of $s_0$. All sums are taken in $\mathbb{R}$ and reduced modulo $2\pi$ to lie again in $\mathbb{T}^{1}$.

%Let us denote $\mu_t:= \text{Law}(x_t)$, the law of $x_t$ on $\mathbb{T}^1$. 
The objective is to control the process in such a way that the distribution $\mu_T$ of the accumulated reward $ \sum_{t=0}^T r_t(s_t,a_t)=s_T$ is as close as possible to a given target distribution $\mu^*$ on $\mathbb{T}^1$. The distance is measured as a distance of characteristic functions. We will be more precise later.
For the purpose of illustration let us assume that $T=1$. In this case we can compute the characteristic function of $s_1$ explicitly. First, consider  $y_1=s_0+a_0+\varepsilon_1.$  Then the law of $y_1$ is a Gaussian centered at $s_0+a_0$ if $a_0$ is deterministic and if $a_0$ is stochastic with distribution $\nu$ then 
$$\PP(y_1\le z)= \int_\RR \Phi(\frac{z-x_0-u}{\sigma}) d \nu(u)  $$
i.e. the distribution of $y_1$ is given by the convolution 
%with $\nu$ and $\delta_{x_0}$ is 
$ \delta_{x_0}*\nu * \Phi.$
Second, we need to consider the nonlinear transformation $w:\RR \to \TT^1$  given by
 $$y\bmod{2\pi}=w(y)= y-2\pi \lfloor \frac{y}{2\pi}\rfloor. $$ 
In such a situation, if  \(y\sim\mathcal{N}(m,\sigma^{2})\) on \(\mathbb{R}\) and 
$s \;=\; w(y)  \;\in\; \mathbb{T}^{1},$ then the density of \(s\) with respect to Lebesgue measure on \([0,2\pi)\) is the
\emph{wrapped Gaussian} (or \emph{circular normal}): 
\[
\Phi^\text{wrap}(x;m,\sigma^{2})
 \;=\;
 \frac{1}{\sqrt{2\pi\sigma^{2}}}
 \sum_{k\in\mathbb{Z}}
     \exp\!\Bigl(-\frac{(x-m+k)^{2}}{2\sigma^{2}}\Bigr)=\frac{1}{\sqrt{2\pi}}\Bigl\{
 1+2\sum_{k=1}^{\infty}e^{\frac{k^2\sigma^{2}}{2}}\cos\bigl(2k\, (x-m)\bigr)\Bigr\} .
\]
Thus, wrapping the distribution in $\TT^1$ gives 
$$\PP(s_1\le z)= \int_{\TT^1} \Phi^\text{wrap}(\frac{z-x_0-u}{\sigma}) d \nu(u).  $$
When we want to match a certain characteristic function we can do it by choosing $\nu$ accordingly. Note that the characteristic function of a distribution $\mu$ on $\TT^1$ is determined by the sequence $\hat{\mu}(n) :=  \int_{\TT^1} e^{-i n u} \, d\mu(u)$ for $n \in \mathbb{Z}$,
%\[
%\hat{\mu}(n) :=  \int_{\TT^1} e^{-i n u} \, d\mu(u), \quad n \in \mathbb{Z},
%\]
see for example \cite[Chapter 1]{rudin1962fourier}.
Further note that the characteristic function of the wrapped Gaussian density centered at \( s_0 \)  is given by
\[
\widehat{\Phi}^{\text{wrap}}(n; s_0,\sigma^2) = e^{-i n s_0} e^{-\frac{1}{2} \sigma^2 n^2}, \quad n \in \mathbb{Z}.
\]
Thus, identifying the characteristic function of $s_1$ with the one of target distribution $\mu^*$ yields $\mu^*(n)=\widehat \mu_1(n)$ for all $ n\in \ZZ$.
%\begin{equation*}
%	\label{eq: Fourier Matching discrete} \widehat \mu^*(n)=\widehat \mu_1(n) \quad \forall\, n\in \ZZ
%\end{equation*}
Using the convolution property of the characteristic function and solving for the policy that we want to determine yields
$$\widehat \nu(n) = \frac{\widehat\mu^*(n)}{\widehat{\Phi^{\text{wrap}}_{s_0,\sigma}}(n)}=\widehat\mu^*(n) \cdot e^{i n s_0}\; e^{\frac{1}{2} \sigma^2 n^2} \quad \forall\, n\in \ZZ. $$


%\begin{remark}
	It is relatively obvious from the discussion above that when we are considering only deterministic controls we just shift the initial condition without any hope to match any nontrivial target distribution.  
%\end{remark}

\subsection{Extension to Continuous Non Compact Spaces} 
Of course as illustrative and appealing as it is, this example is of limited use and we present it only for pedagogical reasons. The main question is \emph{how to extend to continuous non compact Borel spaces with non linear dynamics?} When we cannot trace the functional form of the distribution of controls since we don't have anymore convolutions and the Fourier representation cannot be broken down analytically.  

To motivate our choices in Section \ref{sec:Problem+Ass} we first notice that to extend the problem in case of continuous non compact spaces we need to define a distance between the characteristic functions of the reward $R_T=\sum_{k=0}^{T-1} r_k(s_k,a_k)$ and a target characteristic function $\varphi^*$: 
$$ \int_{\mathbb{R}} \big|\varphi^*(u) - \EE[e^{iuR_T}]\big|^2 \, w(u)\, du,$$
with some weight function $w$ (e.g.\ Gaussian) to make the integral finite and emphasize relevant frequencies. A natural choice for working directly with probability measures could also be the Wasserstein distance but with a serious caveat that it is not differentiable out of the box and thus not easy to use as a loss function for our purposes. The theoretical foundations of problems like this have been laid in \cite{bauerle2025distributionalbellmanequation}. The authors there give a dynamic programming equation for the solution. However, the question still remains how these problems can numerically be solved efficiently.
%The second idea comes to counter our lack of explicit convolutions for tracing the functional form of the distribution of controls. Instead 
Thus, in this paper we use artificial neural networks (ANNs) to parametrize the controls where we condition them on the current state and reward and using some external noise in order to learn to sample optimally from the unknown action distribution $\nu$ when we have no access and or assumption on it. The ANN parameters are trained by a gradient descent scheme. 

 To wrap up, we design a \emph{model free} policy gradient algorithm using only samples from the dynamics and rewards to minimize a characteristic function loss.  

\subsection{Choosing the Target}
There are several ways to choose the target distribution. This may either be a desired risk-return distribution when we think about investment problems where a certain risk profile is targeted. Another setting could be the problem of imitating others: Suppose you can observe the rewards/outcome of an opponent but you do not know her/his actions. Then you can easily obtain her/his characteristic target function of cumulative rewards and use this to find out a control which is consistent with this behavior. A third application is to design a Markov model such that the behavior in terms of cumulative rewards behaves as observed in nature.

\subsection{Literature Review}
Distributional Reinforcement learning has gained in recent years a lot of attraction with several important contributions. 

One of the most influential advances is the quantile-regression formulation of  \cite{dabney2018distributional}, which replaces the fixed-support categorical approximation of C51 from the seminal paper \cite{bellemare2017distributional} with a trainable quantile parameterization. This representation yields a theoretically consistent approximation of value-return distributions in the 1-Wasserstein metric and resolves several limitations of earlier approaches, most notably by providing unbiased stochastic gradient estimates via asymmetric quantile losses and by ensuring that the projected Bellman operator is a contraction.  %In this sense, quantile-regression DRL provides important conceptual background demonstrating that distributional representations and unbiased stochastic gradients can substantially enhance stability yet it does not address the distribution-matching control problem considered here. 

Further along this line of research on \emph{value-based control} the article of \cite{bellemare2020distributional} introduces the first convergence guarantees for a distributional RL algorithm combined with function approximation, addressing a central theoretical gap left open by earlier works such as C51 and its analysis by \cite{bellemare2017distributional}. Its key contribution is the construction of a fully Cram\'er-based loss free of softmax and KL divergence which extends the Cram\'er distance to arbitrary real vectors and augments it with a normalization penalty, enabling a new algorithm (S51) whose updates are mathematically tractable. Using this generalized loss, the authors prove that the projected distributional Bellman operator converges under linear function approximation and quantify the resulting approximation error, showing, perhaps unexpectedly, that distributional methods can yield worse expected-value accuracy than classical TD under function approximation. 

In \cite{achab2023onestepdistributionalreinforcementlearning}, the authors introduce a simplified distributional framework that resolves one of the central theoretical difficulties of classical distributional RL: the instability of the optimality operator in the control case. By restricting attention to the randomness induced by the first transition only, the authors construct one-step distributional Bellman operators that unlike the original distributional operators are $\gamma$-contractions in all Wasserstein metrics for both policy evaluation and control, leading to a unified almost-sure convergence theory even without assuming uniqueness of the optimal policy. Using these operators, they develop new Cram\'er-projected algorithms (tabular and deep variants such as OS-C51) whose categorical updates are simpler, cheaper, and provably convergent.

Nevertheless, this line of work remains fundamentally oriented toward \emph{value-based control}: the goal is to estimate the distribution of the return induced by a fixed policy and to extract improved greedy or risk-sensitive actions from it. By contrast, the problem studied in this article is of a different nature: instead of estimating an endogenous return distribution, we aim to \emph{steer} the distribution of a terminal random variable toward a prescribed target law by directly optimizing the policy. Whereas quantile-based DRL methods rely on Bellman recursions, additive rewards, and the value-function framework, our approach operates on a lifted state space of probability measures and optimizes a characteristic-function divergence akin to a policy iteration method. We show in Example \ref{se:forzen_lake} that our framework is general and flexible enough that it can accommodate solutions to some classical RL problems with minimal modifications. 

Recent work has extended distributional reinforcement learning to the continuous-time setting by characterizing the evolution of return distributions through a distributional analogue of the backward Kolmogorov equation. In \cite{wiltzer2022distributional} the authors derived what they call the distributional HJB equation, but which is in fact a linear backward Kolmogorov PDE for the state-conditioned return CDF under a fixed policy, thereby providing the first continuous-time description of return distributions beyond expectations. To make this infinite-dimensional PDE numerically tractable, they introduced the statistical HJB (SHJB) loss, which replaces the full distribution by a finite set of statistics (e.g., quantiles) and measures the discrepancy between the approximate CDF and the PDE operator. However, that work did not establish whether minimizing the SHJB loss actually produces consistent approximations of the true return distribution. This gap is filled by \cite{alhosh2025tractablerepresentationsconvergentapproximation}, who prove that if the chosen imputation strategy (mapping finite statistics to a distribution) satisfies a mild topological condition, namely convergence in the sense of tempered distributions then SHJB minimization is provably convergent: as the number of statistics increases, the approximate solutions converge to the true solution of the Kolmogorov PDE. Moreover, they show that the widely used quantile representation satisfies this condition and achieves an $\mathcal{O}(\frac{1}{N})$ approximation error for the CDF, thereby giving the first rigorous justification that quantile-based continuous-time DRL is a sound and convergent discretization scheme for the distributional Kolmogorov equation. 

In contrast to these works, which focus primarily on \emph{estimating} the return distribution under a policy (or its continuous-time evolution), \cite{bauerle2025distributionalbellmanequation,pires2025optimizing}  introduce distributional dynamic programming, a general DP framework for \emph{optimizing arbitrary statistical functionals of the 
return distribution}. Their approach is the closest to ours in this article. We agree on the definition of the distributional MDP which augments states with cumulative reward (stock) and combines with a distributional Bellman recursion defined over this augmented MDP, enabling principled optimization of objectives such as 
quantiles, CVaR, and other risk-sensitive functionals. In this way, distributional DP broadens the scope of distributional RL by integrating distributional objectives directly into the dynamic programming principle, complementing earlier work on distributional estimation in both discrete and continuous time. In contrast with \cite{pires2025optimizing}, like in \cite{bauerle2025distributionalbellmanequation} we resort to a lifted distributional MDP (see Section \ref{sec:liftedMDP})  to formally describe the Bellman recursion and identify the class of optimal policies, mainly for two reasons, first one being the deterministic nature of the lifted MDP and the time homogeneous Markovian structure of the augmented MDP.  

Last but not least we mention two works that connect stochastic control with optimal transport and are very close in spirit to ours but different in methods. In \cite{Alouadi2025SBTS} the authors study time series generation through the Schrödinger Bridge (SB) framework, which formulates generative modeling as a stochastic control problem on path space: one selects a drift process minimizing a quadratic control (relative-entropy) cost with respect to a Brownian reference measure, subject to matching prescribed marginal distributions of the state process at given time points. The main conceptual difference lies in the location of the distributional objective: while the SB approach enforces distributional constraints directly on the state process itself, we consider controlled dynamics augmented with a cumulative reward variable and impose the distribution-matching objective on this derived random quantity rather than on the state trajectory. In this sense, rewards act as an additional layer on top of the dynamics, transforming path distributions into outcome distributions. Unlike the SB formulation, we are not aware of any optimal transport or entropic regularization interpretation on path space for our loss. Nevertheless, \cite{Alouadi2025SBTS} is very interesting in conjunction with \cite{wiltzer2022distributional} for a general bridge interpretation. 

 In \cite{TerpinLanzettiDorflers2024DPOT} the authors consider fleets of identical agents whose dynamics are non-interacting, each agent evolves independently according to a common controlled transition map, and the population state evolves as the pushforward of the current distribution through this single-agent dynamics. Any coupling between agents arises exclusively through the objective function, via optimal transport discrepancies that compare the induced population distribution to a prescribed reference measure. As a result, the separation principle identified in that work critically relies on this absence of interaction in the dynamics. Here, we stress once more the difference that we operate in a lifted distributional MDP and rely on a direct policy method. However we find the connection with optimal transport very interesting at least in the case of value iteration. 

 
\subsection{Summary and Organisation of the Article}

The remainder of the article is organized as follows. In Section~\ref{sec:Problem+Ass}, we formally introduce the distributional control problem, define the characteristic–function matching objective, and reformulate it as a lifted distributional Markov Decision Process on an augmented state space. We also specify the class of randomized Markov policies considered and state the standing assumptions on the dynamics, rewards, and neural-network parametrization. Section 3 presents the policy-gradient algorithm for distribution matching, including the discretization of the Fourier domain, the construction of an unbiased (up to finite-sample effects) stochastic gradient estimator, and the full learning procedure. In Section 4, we establish our main theoretical result, proving convergence of the proposed algorithm to a stationary point under standard stochastic approximation conditions. Section~\ref{sec:appl} illustrates the methodology on a series of numerical examples, namely linear–quadratic control, an investment problem, a compactly supported distribution and classical benchmark MDPs, highlighting both the flexibility and the intrinsic non-uniqueness of distributional optimal controls. Technical proofs and auxiliary results are collected in the Appendix.

\section{The Problem and Assumptions}\label{sec:Problem+Ass}
\subsection{Problem Formulation}
The main purpose of this section is to give a formal definition to our general problem. To this end, we assume a finite time horizon $T\in\mathbb{N}$ and let us denote the state space and action space  with $\mathcal{S}$ and $\mathcal{A}$, respectively. We assume that both are Borel subsets of $\RR$.   Let $(\Omega, \mathcal{F}, \mathbb{P})$ be a complete probability space and $(\varepsilon_t)_{t=0}^{T-1}$ a sequence of i.i.d. standard normal random variables defined on this space. For ease of notation we consider a stationary model here, but the extension to non-stationary systems is immediate.
The system evolves according to the stochastic dynamics:
\begin{equation}
\label{eq:def:nonlinear_dynamics}
    s_{t+1} = F(s_t, a_t, \varepsilon_{t+1}), \quad \varepsilon_{t+1} \sim \mathcal{N}(0,1),
\end{equation}
with initial state $s_0$  deterministic or random $\mathcal{F}_0$-measurable. Furthermore, $F:\cS \times \cA \times \RR \to \cS$ is a measurable transition function and $a_{t}\in\mathcal A$ is the action chosen at time point $t$. We assume that the decision maker can decide about   a  policy $\sigma=(\sigma_0,\ldots,\sigma_{T-1})$ where $\sigma_t$ is the randomized, history dependent decision rule at time $t$, i.e. $a_t\sim \sigma_t(\cdot|s_0,a_0,\ldots,s_t)$.  

Assume a measurable one-stage reward $r:\cS \times \cA \to \RR$. The cumulative reward for a trajectory up to time $t$ is:   

\begin{equation}
    R_t = \sum_{k=0}^{t-1} r(s_k, a_k). 
\end{equation}

For this work we choose to see $R_T$ as a random variable and we are interested in its distribution. The distribution depends on the chosen policy $\sigma$, hence we denote it by $\text{Law}^\sigma(R_T)$. 

When we are given a target distribution $\mu^\star$ on $\RR$, our goal is to find a policy $\sigma$  such that 
\begin{equation}
\label{eq:matching equation}
    \text{Law}^\sigma(R_T) \approx \mu^\star.
\end{equation}

To formalize the matching, we introduce a divergence between probability distributions, the weighted squared $L^2$ distance between the characteristic functions:
\begin{equation}
\label{eq:def: general Loss} 
    \mathcal{L}(\sigma) := \int_{\mathbb{R}} \left| \varphi^\star(u) - \varphi_{\sigma}(u) \right|^2 w(u) \, du,
\end{equation}
where $\varphi_{\sigma}(u) := \mathbb{E}^\sigma[e^{iu R_T}]$ is the characteristic function of the real-valued random variable $R_T$ when policy $\sigma$ is chosen, $\varphi^\star$ is the characteristic function of $\mu^*$ and $w: \mathbb{R} \to \mathbb{R}_+$ is a suitable weighting function satisfying the weight moment assumptions 
\begin{equation*}
  \int_\RR  w(u) du < \infty, \quad   \int_\RR |u| \cdot w(u) du < \infty, \quad  \int_\RR u^2 \cdot w(u) du < \infty.
\end{equation*}
Distances like these are popular in statistics for testing distribution hypothesis, \cite{epps1983test,baringhaus1988consistent}.
The problem becomes:
\begin{equation}
\label{eq:general learning problem}
    \min_{\sigma} \mathcal{L}(\sigma).
\end{equation}

Now, we reformulate Problem \eqref{eq:general learning problem} according to \cite{bauerle2025distributionalbellmanequation} as a \emph{lifted MDP} in order to distinguish classes of controls among which we will look for optimal ones. 

 \subsection{Lifted Distributional MDP Formulation} \label{sec:liftedMDP}
 We now recast optimization problem \eqref{eq:general learning problem} as a Markov Decision Process on an \emph{augmented state space} that captures the distribution of outcomes. 

To formulate first the \emph{distributional MDP}, we retain the original notation and introduce an \emph{augmented state variable} $(s_t, R_t)$ taking values in the product space $\mathcal{S} \times \mathcal{R}$, where $\mathcal{R} \subseteq \mathbb{R}$ is the set of possible cumulative rewards. The law of $(s_t, R_t)$ under a policy $\upsigma$ (and fixed initial distribution) is denoted $\mathsf F_t^\upsigma \in \mathcal{P}(\mathcal{S} \times \mathcal{R})$, i.e., the joint distribution of $(s_t, R_t)$. In particular, $\mathsf F_0 = \delta_{s_0 \times 0}$ is the initial distribution of $(s_0, R_0)$. The terminal law $\mathsf F_T^\sigma$ encodes the distribution of both the terminal state $s_T$ and the cumulative reward $R_T$. Note that the marginal of $\mathsf F_T^\sigma$ on the $R$-coordinate is precisely $\text{Law}(R_T)$. 

We can now define a \emph{lifted distributional MDP} whose state at each time is the joint distribution $\mathsf F_t^\sigma$ of  $(s_t, R_t)$. The lifted model is specified as follows:

\begin{enumerate}
\item[a)] \textbf{State space:} $\mathcal{X} := \mathcal{P}(\mathcal{S} \times \mathcal{R})$. An element $\mathsf F \in \mathcal{X}$ is a joint distribution of the state and accumulated reward. At time $t$, the lifted state is $\mathsf F_t = \text{Law}^\sigma(s_t, R_t)$ under the current policy $\sigma$. This lifted state encodes all the necessary information of the process.

\item[b)] \textbf{Action space:} $\Pi^M := \{\pi: \mathcal{S} \times \mathcal{R} \to \mathcal{P}(\mathcal{A})\}$, the set of Markov (possibly randomized) decision rules on the augmented state. That is, $\pi \in \Pi^M$ is a conditional distribution $\pi(da|s, R)$ prescribing an action distribution given state $s$ and cumulative reward $R$. A lifted policy is a sequence $\pi_0, \pi_1, \ldots, \pi_{T-1}$ with $\pi_t \in \Pi^M$.

\item[c)] \textbf{One-stage reward:} $0$ for each $t < T$. Intermediate transitions are assigned zero reward, as only the terminal outcome matters.

\item[d)] \textbf{Terminal reward functional:} $H: \mathcal{P}(\mathcal{S} \times \mathcal{R}) \to \mathbb{R}$ is defined by
\[
H(\mathsf F) := \int_\RR \Big|\varphi^*(u)- \int e^{iux} \mathsf F(\mathcal{S},dx)\Big|w(u)du.
\]

\item[e)] \textbf{State-transition operator:} $T: \mathcal{P}(\mathcal{S} \times \mathcal{R}) \times \Pi^M \to \mathcal{P}(\mathcal{S} \times \mathcal{R})$ describes the evolution of the distribution under a policy. Given $\mathsf F \in \mathcal{P}(\mathcal{S} \times \mathcal{R})$ and $\pi \in \Pi^M$, the next distribution $\mathsf F' = T^{\pi}(\mathsf F)$ is the law of $(s_{t+1}, R_{t+1})$ where:
\[
\begin{aligned}
(s_t, R_t) &\sim \mathsf F, \\
a_t &\sim \pi(\cdot | s_t, R_t), \\
s_{t+1} &= F(s_t, a_t, \varepsilon_{t+1}), \\
R_{t+1} &= R_t + r(s_t, a_t).
\end{aligned}
\]
Formally, for any measurable set $B \subseteq \mathcal{S} \times \mathcal{R}$,
\[
T^{\pi}(\mathsf F)(B) = \int_{\mathcal{S} \times \mathcal{R}} \int_{\mathcal{A}} \mathbb{P}\left\{ (F(s,a,\varepsilon),\; R + r(s,a)) \in B \right\} \pi(da|s,R)\,F(ds,dR).
\]
\end{enumerate}

This defines a finite-horizon deterministic dynamic program on the space of probability measures, with state variable $\mathsf F_t \in \mathcal{P}(\mathcal{S} \times \mathcal{R})$. Any policy $\upsigma$ in the original model induces a sequence $\mathsf F_0 \xrightarrow{\sigma_0} \mathsf F_1 \xrightarrow{\sigma_1} \dots \xrightarrow{\sigma_{T-1}} \mathsf F_T$ with terminal objective $H(F_T)$ and similar any sequence $(\pi_0, \ldots, \pi_{T-1})$ of kernels in the lifted MDP defines a feasible path $\mathsf F_0 \xrightarrow{\pi_0} \mathsf F_1 \xrightarrow{\pi_1} \cdots \xrightarrow{\pi_{T-1}} \mathsf F_T$. It is possible to construct the sequence  $(\pi_0, \ldots, \pi_{T-1})$ such that distribution sequences coincide (see \cite{bauerle2025distributionalbellmanequation} Prop. 2.2)
Hence, the original problem \eqref{eq:general learning problem} is equivalent to:
\begin{equation} \label{eq:equivalent_value_function}
\min_{(\pi_0,\ldots,\pi_{T-1})} H\left(T^{\pi_{T-1}} \circ \cdots \circ T^{\pi_0}(\mathsf F_0)\right).
\end{equation}

For completeness we define the value functional $V_t: \mathcal{P}(\mathcal{S} \times \mathcal{R}) \to \mathbb{R}$ by:
\begin{equation}
	\label{eq:value_function}
V_T(\mathsf F) := H(\mathsf F), \quad
V_t(\mathsf F) := \inf_{\pi \in \Pi^M} V_{t+1}\left(T^{\pi}(\mathsf F)\right), \quad t = T-1, \dots, 0.	
\end{equation}

Then $V_0(\mathsf F_0)$ gives the optimal value of the lifted problem and equals the value in \eqref{eq:general learning problem}.
Finally, a key consequence of this formulation is that the search for an optimal policy may be restricted \emph{without loss of generality} to actions $\pi_t \in \Pi^M$ on  $(s_t, R_t)$. Any history-dependent policy can be equivalently represented as a sequence $(\pi_0, \dots, \pi_{T-1})$ producing the same law of outcomes, see \cite[Remark 2.2]{bauerle2025distributionalbellmanequation}.  
However, due to the complicated state and action space in the lifted MDP it is computationally very hard to perform the value iteration algorithm \eqref{eq:value_function}. In order to circumvent this problem we define a parametrized policy class and use gradient descent to obtain the optimal policy.

\subsection{Controls}
Now, in contrast to value iteration methods based on $V_T(\mathsf F)$ we chose to use $H(\mathsf F)$ as a target function for a parametric policy iteration in line with Motivation \ref{se:motivation}. As we already explained, the lack of analytical traceability for the distribution of the actions calls us for function approximation. 
We will choose artificial neural nets (ANNs) to approximate the decision rules. Their use, in our work is very close in the spirit of \cite{Goodfellow2014GAN}, in particular, we learn an optimal sampling mechanism from the unknown distribution of controls, while making no assumption on their distribution. We stress that our method is likelihood and model free. Very relevant to ours but with a direct paramterization is \cite{han2016deep,hure2021deep}. In our work, we make use of Lipschitz assumptions and properties of the MDP to prove convergence of the algorithm. This has also been done in \cite{pirotta2015policy} where, however the authors parametrize the probability law of the action. In \cite{fatkhullin2023stochasticpolicygradientmethods}, the authors deal as well with an implicit method but assumpting a density.

More precisely  we further assume that on our  probability space $(\Omega, \mathcal{F},\mathbb{P})$ there exists another  sequence of i.i.d.\  standard normal random variables $(z_t)$ which is independent from $(\varepsilon_t)$.
We define the control $a_t$ at time $t$, sampled from the Markovian policy $\pi_t(\cdot|s_t,R_t)$  as 
\begin{equation}
    a_t = f(\theta,s_t,R_t,z_t,t), %\quad z_t \sim \mathcal{N}(0,1), 
\end{equation}
where $f: \Theta\times \mathcal{S}\times \mathcal R \times \mathbb{R} \times \{0,1,\ldots,T-1\}\to \mathcal{A}$ is a neural network parametrised by $\theta\in\Theta$. Note that we need $z_t$ to generated the randomness. W.l.o.g.\ we could also use a random variable uniformly distributed over $(0,1)$ or any other choice will do. 

To summarize we have the following evolution for $t=0,1,\ldots ,T-1$ of the parametrized system where $s_0$ is a given state and $R_0=0$:
\begin{eqnarray*}
a_t^\theta &=& f(\theta,s_t^\theta,R_t^\theta,z_t,t),\\
s_{t+1}^\theta &=& F(s_t^\theta,a_t^\theta, \varepsilon_{t+1}),
\end{eqnarray*}
and for $t=1,\ldots ,T$
$$R_t^\theta = \sum_{k=0}^{t-1} r(s_k^\theta,a_k^\theta).$$ 

In general, we consider a feedforward neural network with $K$ hidden layers of width $P$, taking as input $(s,R,z,t) \in \mathcal{S}\times \mathcal R \times \mathbb{R} \times \{0,1,\ldots,T-1\}$,\footnote{in practice we can either take one "big" neural network including time as a separate input variable or have $T$ different neural networks each with a set of parameters $\theta_t$, the two specifications give identical results} applying the activation function component-wise and producing a scalar output.
While our results hold for arbitrary $K, P$ for the sake of readability we assume a 2-layer neural network  

\begin{equation}
	\label{eq:definition_NN_formula}
	f(\theta,s,R,z,t) =  w^2\sigma(w^1_s s + w^1_R R+ w^1_z z + w^1_t t+ b^1)+b^2. 
\end{equation}
where $\theta=(w^1_s,w_R^1,w_z^1,w_t^1,w^2,b^1,b^2)\in \RR^7.$ For convenience we will write
\begin{equation}
	\label{eq:definition_NN_formula2}
	f(\theta,x) =  w^2\sigma(W^1 x + b^1)+b^2 
\end{equation}
for $W^1=(w^1_s,w_R^1,w_z^1,w_t^1)\in \RR^4$ and $x=(s,R,z,t)\in \mathcal{S}\times \mathcal R \times \mathbb{R} \times \{0,1,\ldots,T-1\}.$
 The convergence result for the algorithm that we later show also holds for multi-layer networks. We only need to readjust the constants to transfer the regularity via the different layers of the ANN. 

\subsection{Assumptions and Some Implications}
Throughout we need several assumptions on the model data such that we can prove convergence of our algorithm.
\subsubsection{Assumptions on the neural net data}
We assume that the parameters $\theta$ of the neural net are taken from a compact set $\Theta\subset \RR^7$ and that the activation function $\sigma$ shows a high regularity in the sense that the following assumption holds.

\begin{assumption}
	\label{as:1} The parameter set $\Theta$ is compact and $\sigma, \sigma'$ are Lipschitz-continuous and non-decreasing. We allow that $\sigma'$ does not exist at isolated points.
\end{assumption}

Note that for example the activation function $\sigma(x)=\tanh(x)$ or ReLu satisfy the requirements. 
%Moreover, all actions which can be generated from \eqref{eq:definition_NN_formula2} are always in a compact subset of $\RR$.

\subsubsection{Assumptions on the growth of the state process}
We do not want to bound the state process because this would rule out a number of interesting applications. Instead we assume that the Markov Decision Process possesses a bounding function. A concept which is often used to bound the value function (see e.g. \cite{BauerleRieder}, Sec. 7.3).

\begin{assumption}
	\label{as:2} There exists a linear bounding function  for the MDP, i.e.\ there exist $c_b,d_b\in\RR_+$ such that for all $s\in \mathcal{S}$ and $\varepsilon \sim \cN(0,1)$
%    \begin{itemize}
  %      \item[(i)] $|s| \le c_s b(s),$
 %        \item[(i)] $|r(s,a)| \le c_r |s|$ for all $a\in \mathcal{A}$,
 %         \item[(ii)] 
 $$\EE\left[|F(s,a,\varepsilon)| \right] \le c_b |s| + d_b \mbox{ for all } a\in \mathcal{A}$$
%    \end{itemize}
\end{assumption}

From the assumption of a bounding function we now obtain:
\begin{eqnarray*}
    \EE[|s_t|] &=& \EE\left[ |F(s_{t-1},a_{t-1},\varepsilon_{t})| \right] \\
    &=&  \EE\big[\EE\left[ |F(s_{t-1},a_{t-1},\varepsilon_t)|| s_{t-1},a_{t-1}\right]\big]\\
    &\le & c_b \EE[|s_{t-1}|] + d_b\le \ldots \le  c_b^t |s_0|+ d_b \sum_{k=0}^{t-1}c_b^k.
\end{eqnarray*}

In the next subsection we state a number of Lipschitz-continuity properties of our data. 

\subsubsection{Further Lipschitz-Continuity Assumptions on the Data}
For the convergence of the algorithm we have to impose further Lipschitz-continuity properties of the data.

\begin{assumption}
	\label{as:3} The transition function $F:\cS \times \cA \times \RR \to \cS$ is Lipschitz-continuous in the third component, uniformly in $s,a,$ i.e. there exits a constant $L_F>0$ such that
    $$ |F(s,a,\varepsilon)-F(s,a,\tilde\varepsilon)| \le L_F |\varepsilon-\tilde\varepsilon|, \quad \mbox{for all } s\in \cS,  a\in \cA, \varepsilon, \tilde\varepsilon \in\RR.$$
\end{assumption}

\begin{assumption} \label{ass:rFLip}
We assume
\begin{itemize}
\item[(i)]$r$ is Lipschitz continuous, i.e. for all $s,\tilde s\in \cS,a,\tilde a \in \cA $ there is a $L_r>0$ such that
$$|r(s,a)-r(\tilde s,\tilde a)| \le L_{r}\big( |s-\tilde s| + |a-\tilde a| \big) $$
\item[(ii)] $F$ is  Lipschitz continuous, i.e. for all $s,\tilde s\in \cS,a, \tilde a \in \cA $ and $\varepsilon \in\RR$ there is a $\tilde L_F>0$ such that
$$|F(s,a,\varepsilon)-F(\tilde s,\tilde a,\varepsilon)| \le \tilde L_{F} \big( |s-\tilde s| + |a-\tilde a|\big). $$
\end{itemize}
\end{assumption}

We denote by $\partial_s r, \partial_a r$ and $\partial_s  F, \partial_a F $ the derivative w.r.t.\ the corresponding component and assume the following:

\begin{assumption} \label{ass:rFLipgrad}
We assume
\begin{itemize}
\item[(i)]$\partial_s r, \partial_a r$ exist and are  Lipschitz continuous, i.e. for all $s,\tilde s\in \cS,a, \tilde a \in \cA $ and for $i=s,a$
$$|\partial_i r(s,a)-\partial_ir(\tilde s,\tilde a)| \le L_{\nabla}^{r,i} \big( |s-\tilde s| + |a-\tilde a|  \big)$$
\item[(ii)] $\partial_s  F, \partial_a F$ exist and are Lipschitz continuous, i.e. for all $s,\tilde s\in \cS,a, \tilde a \in \cA,\varepsilon,\tilde\varepsilon\in\RR $ for $ i=s,a$
$$|\partial_iF(s,a,\varepsilon)-\partial_i F(\tilde s,\tilde a,\tilde\varepsilon)| \le L_{\nabla}^{F,i} \big( |s-\tilde s| + |a-\tilde a|\big). $$
\end{itemize}
\end{assumption}

\begin{remark}
    First note that if state and action space are compact, Assumption \ref{as:2} would be satisfied. Further the condition that reward $r$ and transition function $F$
 are continuously differentiable would directly imply Assumptions \ref{ass:rFLip} and \ref{ass:rFLipgrad}.
 \end{remark}

\subsubsection{Bounds on the evolution of the process}
From what we assumed so far, we can get bounds on the state evolution. More precisely, using the Gaussian Concentration Inequality (see Appendix \ref{sec:concentration}) we obtain for constants $\eta_{t+1}>0:$
\begin{eqnarray*}
   \PP (|s_{t}-\EE[s_{t}]| > \eta_{t})  &=&  \PP (|F(s_{t-1},a_{t-1},\varepsilon_{t})-\EE[F(s_{t-1},a_{t-1},\varepsilon_{t})]| > \eta_{t})\\
   &=&  \EE\Big[  \PP (|F(s_{t-1},a_{t-1},\varepsilon_{t})-\EE[F(s_{t-1},a_{t-1},\varepsilon_{t})]| > \eta_{t}\; |\; s_{t-1},a_{t-1}) \Big]\\
   &\le & 2 \exp\!\Big(-\frac{\eta_{t}^2}{2\,L_F^2}\Big).
\end{eqnarray*}
Suppose $\gamma_{t}>0$ is given. If we choose $\eta_t = L_F \sqrt{ 2 \ln(2/\gamma_t)}$, then $ \PP (|s_{t}-\EE[s_{t}]| > \eta_{t})\le \gamma_t.$
Next let us define the event that the state process is not too far away from its expectation
\begin{equation}\label{eq:goodevent}
    G' := \bigcap_{t=1}^T \Big\{|s_t-\EE[s_t]| \le  \eta_t \Big\}
\end{equation} 
We obtain
\begin{eqnarray*}
    \PP(G') &=& 1- \PP\Big( \bigcup_{t=1}^T \{|s_t-\EE[s_t]| >  \eta_t \}\Big) \\
    &\ge& 1- \sum_{t=1}^T \PP\big(  \{|s_t-\EE[s_t]| >  \eta_t \}\big) \ge 1- \sum_{t=1}^T  \gamma_t%2 \sum_{t=1}^T  \exp\!\Big(-\frac{\eta_t^2}{2\,L_F^2}\Big).
\end{eqnarray*}
%Choosing $(\eta_t)$ large enough we can make the probability of $G$ as large as we like. 
Further we bound the action process by choosing $B^z>0$ and define
$$ G:= G' \cap \bigcap_{t=1}^T \Big\{|z_t| \le  B^z \Big\}$$
In what follows we assume that $\PP(G)\ge 1-\gamma $ for a given $\gamma>0$. This can be achieved by choosing $B^z$ and $\eta_t$ large enough.
%with $\gamma=\sum_{t=1}^T  \gamma_t$. 
%As long as we consider realizations in $G$, also the reward process $(R_t)$ is bounded.

 Instead of assuming a compact state  space we work with what we call the good event $G$. As long as $\omega\in G$ we have a bounded state and action process  and thus also bounded rewards. More precisely we obtain on $G$:
 \begin{equation}\label{eq:s_bound}
     |s_t| \le \EE[|s_t|] + L_F \sqrt{ 2 \ln(2/\gamma_t)}\le c_b^t |s_0|+ d_b \sum_{k=0}^{t-1}c_b^k+ L_F \sqrt{ 2 \ln(2/\gamma_t)} =: B_t^s
 \end{equation} 

\section{Gradient Descent Algorithm}
In order to define a gradient descent scheme for the parameters of the ANN, we first reformulate the objective \eqref{eq:general learning problem} given our ANN parametrisation as  

\begin{equation}
	\label{eq: parametrised Loss}
	\min_{\theta \in \RR^7} \cL(\theta) = \min_{\theta \in \RR^7} \int_{\mathbb{R}} \left| \varphi^\star(u) - \EE[e^{iu R_T^\theta}]  \right|^2 w(u) \, du.
\end{equation}
For the gradient descent we need the gradient of $\cL$ with respect to the parameters of the ANN. It can be computed by a pathwise derivative:
\begin{align}
\label{eq:def:Gradient_Loss}
   \nabla_\theta \mathcal{L}(\theta) =
   2 \int_{\mathbb{R}} \operatorname{Re} \left[ \overline{\left( \varphi^\star(u)-\EE[e^{iu R_T^\theta}]   \right)}\cdot \mathbb{E}\left[i u e^{i u R_T^\theta} \cdot \nabla_\theta R_T^\theta \right] \right] w(u) du
\end{align}

In order to proceed we use three approximations for the gradient: 
\subsection{Conditioning on the Good Event}
We first condition the expectation on the good event $G,$ i.e.\ instead of $\EE[e^{iu R_T^\theta}]$ etc.\ we consider
$$ \varphi_\theta(u):=\EE_G[e^{iu R_T^\theta}]:= \EE[e^{iu R_T^\theta}|G].$$
%In what follows we denote by $\cL_G$ the objective function  where we replace $\EE[\cdot]$ by $\EE_G[\cdot]$ in \ref{eq: parametrised Loss}.
We are able to control the error that we make by this conditioning irrespective of $\theta.$ We discuss this in more detail in the next remark.

\begin{remark}[Error induced by conditioning on $G$]
Suppose $G_n$ is an increasing sequence of subset of $\Omega$ such that $G_n \uparrow \Omega$ and $\PP(G_n)=1-\gamma_n$ where necessarily $\gamma_n\downarrow 0.$ W.l.o.g.\ we assume $\gamma_0=1/2$. Denote $v_n(\theta) :=  \int_{\mathbb{R}} \left| \varphi^\star(u) - \EE_{G_n}[e^{iu R_T^\theta}]  \right|^2 w(u) \, du$ and  note that  $|e^{iu R_T^\theta}| \le 1$. In Lemma  \ref{lem:LR} we show that $\theta\mapsto R_T^\theta$ is continuous, thus due to dominated convergence $\theta\mapsto  v_n(\theta)$ is also continuous. Hence minimum points $\theta_n^*$ of $v_n$ exist for all $n\in\NN$ because $\Theta$ is compact. Moreover, we obtain for $n\ge m$ and $\theta\in \Theta$:
\begin{eqnarray*}
    |v_n(\theta)-v_m(\theta)| &\le & 4  \int \Big| \EE_{G_n}[e^{iu R_T^\theta}] -  \EE_{G_m}[e^{iu R_T^\theta}]  \Big|  w(u) du \\
    &\le &  16 \gamma_m \int_\RR  w(u) du < \infty
\end{eqnarray*}
  The right hand side obviously does not depend on $\theta$ and converges to zero for $m\to\infty.$ Thus, Thm A.1.5 in \cite{BauerleRieder} implies that any sequence of minimum points $(\theta_n^*)$ has converging subsequences with limits being optimal for $v(\theta)= \cL(\theta).$  
\end{remark}

%Further we approximate in the general setting (in special cases we obtain explicit formulas) the integral by
%$$ \cL_{G,L}(\theta) = \sum_{\ell=1}^L \beta_\ell  \left| \varphi^\star(u_\ell) - \EE_G[e^{iu_\ell R_T^\theta}]  \right|^2 $$

%In what follows we will approximate the gradient.
%with
%\begin{equation}
%\label{eq:def:path-wise derivative}
%    \nabla_\theta \varphi_{\theta}(u) = \mathbb{E}\left[i u e^{i u R_T^\theta} \cdot \nabla_\theta R_T^\theta \right].
%\end{equation}
%For the next lemma we consider the conditional expectation on $G$, i.e. we write $\EE_G[\cdot] = \EE[\cdot | G]$ and define  
%\begin{align}
%\label{eq:def:Gradient_Loss2}
%    \nabla_\theta \mathcal{L}_{G,L}(\theta) =
%    2 \sum_{\ell=1}^L  \operatorname{Re} \left[ \overline{\left(\EE_G[e^{iu_\ell R_T^\theta}]- \varphi^\star(u_\ell)   \right)}\cdot  \mathbb{E}_G\left[i u_\ell e^{i u_\ell R_T^\theta} \cdot \nabla_\theta R_T^\theta \right] \right].
%\end{align}

\subsection{Discretization of the Integral}
In order to implement the algorithm we have to discretize the integral, unless we are in a special case which we discuss in Remark \ref{rem:Epps} below. We fix $K\in(0,\infty)$ and a quadrature rule on $[-K,K]$
given by nodes $(u_\ell)_{\ell=1}^L$ and nonnegative weights
$( \beta_\ell)_{\ell=1}^L$ such that, for integrable $\mathrm h$,
\[
\int_{-K}^{K} \mathrm h(u)\,w(u)\,du
\;\approx\;
\sum_{\ell=1}^L  \beta_\ell\,\mathrm h(u_\ell).
\]
Here $w(u)$ is the weight from the definition of $\cL$ and $ {\beta}_\ell  $ the  weights. $K$ should be chosen large enough such that the missing integral term is small. Thus, we consider
\begin{equation}\label{eq:LL}
  \cL_{L}(\theta)
:=\sum_{\ell=1}^{L} \beta_\ell\,\big|\varphi^\star(u_\ell)- \varphi_\theta(u_\ell) ]\big|^2.  
\end{equation}

\subsection{Grid-based Stochastic Gradient}
Finally, by coupling $(\varepsilon_{t+1},z_t)$ we draw $M\in\NN$ iid samples   of 
$R_T^{\theta,1},\ldots,R_T^{\theta,M}$ on the set $ G$ and define the estimator for the gradient $\nabla_\theta \mathcal{L}_L(\theta)$ as follows. First define the empirical characteristic function and its derivative by
$$ \hat\varphi_\theta(u) = \frac1M \sum_{j=1}^M e^{iu R_T^{\theta,j} }, \quad \widehat{\nabla_\theta \varphi_{\theta}}(u)= \frac1M \sum_{j=1}^M i u e^{iu R_T^{\theta,j}} \cdot \nabla_\theta R_T^{\theta,j}  $$
 Thus, in the end we consider
 \[
\hat\cL_{L}(\theta)
:=\sum_{\ell=1}^{L} \beta_\ell\,\big|\varphi^\star(u_\ell)-\hat\varphi_\theta(u_\ell)\big|^2.
\]
with gradient
 \begin{equation}\label{eq:approxgradient}
   g(\theta)
:= \nabla_\theta \hat\cL_{L}(\theta)= 2\sum_{\ell=1}^{L} \beta_\ell
\mathrm{Re}\Big(
 \overline{(\varphi^*(u_\ell)-\hat\varphi_\theta(u_\ell))} \cdot \widehat{\nabla_\theta \varphi_{\theta}}(u_\ell)
\Big).  
 \end{equation}

\begin{remark}\label{rem:Epps}
A special pleasant case appears when we choose a standard normal as target and weight function. I.e. we have
$$ \varphi^*(u) = e^{-1/2 u^2}, \quad w(u)= 1/\sqrt{2\pi} e^{-1/2 u^2}, \quad u\in\RR.$$
In this case we can skip the  procedure of integral approximation  since the empirical target function has a very explicit representation. More precisely
%, denote by 
%$$ \hat\varphi_\theta(u) = \frac1M \sum_{j=1}^M e^{iu R_T^{\theta,j}}$$
%the empirical characteristic function. Then 
the empirical target function has the following explicit representation (see e.g. \cite{epps1983test,baringhaus1988consistent}):
\begin{eqnarray*}
  && \int_{\mathbb{R}} \left|      \varphi^\star(u)-\hat\varphi_\theta(u) \right|^2 w(u) \, du\\
&=& \frac{1}{M^2} \sum_{j,k=1}^M \exp(-\frac12 |R_T^{\theta,j}-R_T^{\theta,k}|^2) -\sqrt{2}  \frac1M  \sum_{j=1}^M \exp(-1/4 |R_T^{\theta,j}|^4) + \frac{1}{\sqrt{3}}.
\end{eqnarray*} 
%Hence in the algorithm we can replace the empirical gradient by 
%\begin{eqnarray*}
%  g(\theta) 
%&=& \frac{1}{M^2} \sum_{j,k=1}^M \exp(-\frac12 |R_T^{\theta,j}-R_T^{\theta,k}|^2)  (R_T^{\theta,k}-R_T^{\theta,j}) (  \nabla_\theta  R_T^{\theta,j}  -\nabla_\theta R_T^{\theta,k}) \\
%&& +\sqrt{2}  \frac1M  \sum_{j=1}^M \exp(-1/4 |R_T^{\theta,j}|^4) (R_T^{\theta,j})^3 \nabla_\theta  R_T^{\theta,j}. 
%\end{eqnarray*}
\end{remark}
%%%%%%%%%%%
\subsection{Algorithm}
We conclude this section with the complete algorithm for our learning scheme. We want to stress again that the policy defines a sampling procedure, making it a likelihood-free generative model. 


\begin{algorithm}[H]
\DontPrintSemicolon
\caption{Policy Gradient for distribution matching}
\label{alg:theoretical}

\KwIn{Target cumulative rewards  (or target law); policy $a_t=f(\theta,s_t,R_t,z_t,t)$; horizon $T$; learning rate $(\alpha_k)_{k=0}^\infty$;  frequency grid $(u_\ell)_{\ell=1}^{L}$ with weights $\beta_\ell$; batch size $M$.}
\KwOut{Trained parameters $\theta$.}

Compute target empirical CF, $\varphi^*$  from target data or given distribution.\\

Initialize $\theta_0$\;

\For{$k=0,1,2,\dots$}{
  
  \For{$\ell=1$ \KwTo $L$}{
  Simulate $M$ \emph{independent} trajectories for this frequency\\
    \For{$j=1$ \KwTo $M$}{
      $s_0^{\theta_k,j} \leftarrow s_0$, \ $R_0^{\theta_k,j} \leftarrow 0$\;
      \For{$t=0$ \KwTo $T-1$}{
        Sample $z_t^{j} \sim \mathcal{N}(0,1)$,\quad
        $a_t^{\theta_k,j} \leftarrow f\big(\theta_k,s_t^{\theta_k,j},R_t^{\theta_k,j},z_t^{j},t\big)$\;
        Sample $\varepsilon_{t+1}^{j} \sim \mathcal{N}(0,1)$,\quad
        $s_{t+1}^{\theta_k,j} \leftarrow F\!\big(s_t^{\theta_k,j},a_t^{\theta_k,j},\varepsilon_{t+1}^{j}\big)$\;
        $R_{t+1}^{\theta_k,j} \leftarrow R_t^{\theta_k,j} + r\!\big(s_t^{\theta_k,j},a_t^{\theta_k,j}\big)$\;
      }
    }
    
    $\displaystyle \widehat{\varphi}_{\theta_k}(u_\ell) \leftarrow \frac{1}{M}\sum_{j=1}^M e^{\,i\,u_\ell R^{\theta_k,j}_T}$\;

  }
Compute loss:
  \[
  \hat\cL_{L}(\theta_k) = \sum_{\ell=1}^{L} \beta_\ell\,
  \big| \varphi^{*}(u_\ell)-\hat{\varphi}_{\theta_k}(u_\ell) \big|^2
  \]
  Update parameters: $\theta_{k+1} \leftarrow \theta_k - \alpha_k \, g(\theta_k)$\;
  \If{$\hat \cL_L(\theta_k) < \text{threshold} $}{\textbf{break}}
}

\Return $\theta_{k+1}$\;
\end{algorithm}


\begin{remark}
	While details about the specific choice of the input parameters will be given in Section \ref{sec:appl}, we want to comment that alternatives are possible for the implementation and our goal here is not to exhaustively tests all of them. Firstly, the outer look over all nodes $u_\ell$ can be avoided using the same noise per node making the computation cheaper without affecting biasedness of the estimator. Secondly, the variance of the estimator depends on the choice of the weighing function $w(\cdot)$ and the quadrature weights $\beta_\ell$ as it will become apparent in the following section. Of course, we can use in addition, other techniques to couple the noises along iterations to reduce the variance.   
\end{remark}


The main result of our paper is now the convergence of the algorithm against a stationary point. 

\section{Convergence of the Algorithm to a Stationary Point}
In this section we state and prove our main convergence theorem. The problem \eqref{eq: parametrised Loss} is highly non-convex with respect to parameter $\theta$, and thus the best we can hope is convergence to a stationary point. We want also to emphasise that even if the loss was convex this will give a unique global minimum for the loss but no further information on the set of optimal controls. The optimal control doesn't need to be unique as  already noted in \cite[Chapter 7.4, 7.5]{Bellemare2023DistributionalRL} and we will observe this also in Section \ref{sec:appl}. In order to address questions of uniqueness of optimal controls or even define a notion of uniqueness adapted to our needs we should investigate further the value function as defined in \eqref{eq:value_function}, in the spirit of  \cite{bauerle2025distributionalbellmanequation}.\footnote{We leave this for future works} 

Our proof follows the lines of Robbins Monro stochastic approximation result see for instance \cite{Borkar2025StochasticApproxRL} and \cite[Chapter 5 and Exercise 5.30]{Bach2024LearningTheory}. The main difficulty is to prove the Lipschitzness of the gradient of \eqref{eq:LL}. To this end, we need to work in steps, transfering the regularity from our control problem to the characteristic functions and finally to the gradients. One main difficulty, is that the noise, in both control and dynamics (independently) can accumulate in the system and cause divergence. To account for this case, we need to control the noise with the help of the \emph{good event} on which  the convergence happens.       

We first prove the L-smoothness of the target function on $G$, i.e. the fact that the gradient of the target function is Lipschitz-continuous on $G$. The proof consists of a number of lemmas and is deferred to the appendix. In what follows $\|\cdot\|_2$ is the usual Euclidian norm.

\begin{theorem}\label{lem:Lipschitz_gradient}
Under Assumptions 1-5
the mapping $\theta \mapsto \nabla_\theta \mathcal{L}_{L}(\theta)$  is   globally Lipschitz continuous on $\Theta$ with a constant $\mathbf{L}>0$, i.e.
$$\|  \nabla_\theta \mathcal{L}_L(\theta) -  \nabla_\theta \mathcal{L}_L(\tilde\theta)\|_2 \le \mathbf{L} \|\theta-\tilde\theta\|_2 \mbox{ for all }  \theta,\tilde\theta\in\Theta.$$
\end{theorem}

Before we continue we need some estimates. Note that in the algorithm we have to simulate the $z$ and $\varepsilon$ random variables for each iteration step $k$, exactly $M$ times over the planning period $t=0,\ldots,T.$ Thus, we denote these random variables by $z_t^{k,j}$ and $\varepsilon_t^{k,j}$. In what follows we denote by $\cF_k = \sigma(\varepsilon_t^{1,j},\ldots,\varepsilon_t^{k,j},t=0,\ldots, T-1, j=1,\ldots,M) \vee \sigma(z_t^{1,j},\ldots ,z_t^{k,j}, t=1,\ldots,T, j=1,\ldots,M)$ the information available at iteration time $k$ of the algorithm, i.e.\ after $k$ simulations of the full state trajectory in the algorithm  and we denote by $\cF_k^G =\{ G\cap A: A\in \cF_k\}$ the trace $\sigma$-algebra on the good event.

\begin{lemma}\label{lem:bias}
    The bias of the estimator for the gradient of the target function in \eqref{eq:approxgradient} at iteration time $k$ is given by
\begin{align}
\label{eq:bias-decomposition}
\EE\!\big[g(\theta_k)\mid \cF_k^G\big]
=
\nabla_\theta \cL_{L}(\theta_k)
+\mathrm{Bias}_M(\theta_k),
\end{align}
where
\begin{align}
\label{eq:bias-term}
\mathrm{Bias}_M(\theta_k)
=
\frac{2}{M}\sum_{\ell=1}^{L}\beta_\ell\,
\text{Re}\Big(
iu_\ell\,\EE[\nabla_\theta R_T^{\theta_k}|\cF_k^G]
-\overline{\varphi_{\theta_k}(u_\ell)}\,\nabla_\theta \varphi_{\theta_k}(u_\ell)
\Big),
\end{align}
and it satisfies the bound ($B_\nabla^R$ is the constant from Lemma \ref{lem:bound1} )
\begin{equation}
\label{eq:bias-bound}
\big\|\mathrm{Bias}_M(\theta_k)\big\|_2
\;\le\;
\frac{\hat{S}}{M} , \quad \hat{S}:=4\,B_\nabla^R\Big(\sum_{\ell=1}^{L}|\beta_\ell u_\ell|\Big).
\end{equation}
Moreover, we have
\begin{equation}
\label{eq:square-bound} \EE\!\big[ \|g(\theta_k)\|_2^2\mid \cF_k^G\big]\le \hat{S}^2.%16 (B_\nabla^R)^2 \Big(\sum_{\ell =1}^L |\beta_\ell u_\ell|\Big)^2
\end{equation}
\end{lemma}

For the the main convergence result we need some further assumptions on the step sizes (which are the usual ones). 

\begin{assumption} \label{ass:ConvAlgo}
 Let $(\alpha_k), \alpha_k> 0$ be a sequence of step sizes such that $$\sum_{k=0}^\infty \alpha_k = \infty\quad \mbox{ and } \quad \sum_{k=0}^\infty \alpha_k^2 <\infty. $$
\end{assumption}
In what follows we denote $A_K := \sum_{k=0}^{K-1} \alpha_k$ and $\hat{S}$ as in the previous lemma.

\begin{theorem}[Convergence to a stationary point]
\label{thm:conv-discrete} Under Assumptions 1-6 we obtain for the iterated sequence of parameters $(\theta_k)$ from our algorithm given by $\theta_{k+1} \leftarrow \theta_k - \alpha_k \, g(\theta_k)$ that 
\begin{eqnarray}\label{eq:prob-bound-good}
\sum_{k=0}^{K-1}\frac{\alpha_k}{A_K}  \EE_G\Big[\|\nabla_\theta \cL_L(\theta_k)\|_2^2 \Big] &=& \frac{2}{A_K} \cL_L(\theta_0)+ \mathbf{L} \hat{S}^2  \frac{\sum_{k=0}^{K-1} \alpha_k^2}{A_K} + \frac{\hat{S}^2}{M^2}   \\
&=&O\Big(\frac{1}{A_K}\Big)+O\Big(\frac{1}{M^2}\Big)  \mbox{ a.s. for }K,M\to \infty,
\end{eqnarray}
\end{theorem}

Thus, when we increase the number of iterations and the number of samples for the empirical gradient, then the average expected gradient tends to zero and thus the parametrization of the policy via $\theta_k$ to a stationary point. Note that for $\alpha_k=1/k$ we have $A_K=O(\ln(K))$ i.e. for the convergence it is more important to increase the number of iterations than the number of samples per gradient. See \cite{Liu2023HighProbSGD} for a discussion on how to improve convergence rates. \\

\begin{proof}
Fix $k$ and work conditionally on $\cF_k^G$ so that $\theta_k$ is deterministic.
By $\mathbf{L}$--smoothness of $\cL_{L}$,
for $\theta_{k+1}=\theta_k-\alpha_k g(\theta_k)$ we have the standard descent inequality
\begin{equation}
\label{eq:smooth-descent}
\cL_{L}(\theta_{k+1})
\le
\cL_{L}(\theta_k)
-\alpha_k\big\langle\nabla \cL_{L}(\theta_k),\, g(\theta_k)\big\rangle
+\frac{\mathbf{L}}{2}\alpha_k^2\|g(\theta_k)\|_2^2.
\end{equation}
Taking conditional expectation $\EE[\cdot\mid \cF_k^G]$ and using linearity yields
\begin{align}
\label{eq:cond-step}
\EE\!\big[\cL_{L}(\theta_{k+1})\mid \cF_k^G\big]
&\le
\cL_{L}(\theta_k)
-\alpha_k\Big\langle\nabla \cL_{L}(\theta_k),\,\EE[g(\theta_k)\mid \cF_k^G]\Big\rangle
+\frac{\mathbf{L}}{2}\alpha_k^2\EE\!\big[\|g(\theta_k)\|_2^2\mid \cF_k^G\big].
\end{align}

Define the conditional bias vector
\[
b_k
:=
\EE[g(\theta_k)\mid \cF_k^G]-\nabla \cL_{L}(\theta_k).
\]
Then
\[
\Big\langle\nabla \cL_{L}(\theta_k),\,\EE[g(\theta_k)\mid \cF_k^G]\Big\rangle
=
\|\nabla \cL_{L}(\theta_k)\|_2^2
+
\big\langle\nabla \cL_{L}(\theta_k),\, b_k\big\rangle.
\]

Apply Young's inequality $|\langle u,v\rangle|\le \tfrac12\|u\|_2^2+\tfrac12\|v\|_2^2$ to obtain
\[
-\alpha_k\big\langle\nabla \cL_{L}(\theta_k),\, b_k\big\rangle
\le
\frac{\alpha_k}{2}\|\nabla \cL_{L}(\theta_k)\|_2^2
+\frac{\alpha_k}{2}\|b_k\|_2^2.
\]
Plugging this into \eqref{eq:cond-step} gives the one-step inequality
\begin{align}
\label{eq:one-step-final}
\EE\!\big[\cL_{L}(\theta_{k+1})\mid \cF_k^G\big]
&\le
\cL_{L}(\theta_k)
-\frac{\alpha_k}{2}\|\nabla \cL_{L}(\theta_k)\|_2^2
+\frac{\alpha_k}{2}\|b_k\|_2^2
+\frac{\mathbf{L}}{2}\alpha_k^2\EE\!\big[\|g(\theta_k)\|_2^2\mid \cF_k^G\big].
\end{align}

Now use Lemma \ref{lem:bias} to obtain
\[
\frac{\alpha_k}{2}\|b_k\|_2^2
\le
\frac{\alpha_k}{2}\cdot \frac{\hat{S}^2}{M^2}, \quad 
\frac{\mathbf{L}}{2}\alpha_k^2\EE\!\big[\|g(\theta_k)\|_2^2\mid \cF_k^G\big]
\le
\frac{\mathbf{L}}{2} \, \alpha_k^2 \hat{S}^2.
\]
Substitute these into \eqref{eq:one-step-final}, and then take $\EE_G[\cdot]$ to remove the conditioning on $\cF_k$ and rearrange the equation:
\begin{align}
\label{eq:avg-bound-pre}
\frac{\alpha_k}{2}\EE_G\!\Big[\|\nabla \cL_{L}(\theta_k)\|_2^2\ \Big]
&\le
\EE_G\!\Big[\cL_{L}(\theta_k)-\cL_{L}(\theta_{k+1})\ \Big]
+\frac{\mathbf{L}}{2}  \, \alpha_k^2 \hat{S}^2
+\frac{\alpha_k}{2}\cdot \frac{\hat{S}^2}{M^2}.
\end{align}

Sum \eqref{eq:avg-bound-pre} from $k=0$ to $K-1$. The loss term telescopes and $\cL_{L}\ge 0$ gives
\begin{align}
\label{eq:sum-bound}
\frac12\sum_{k=0}^{K-1}\alpha_k\EE_G\!\Big[\|\nabla \cL_{L}(\theta_k)\|_2^2\ \Big]
&\le
\cL_{L}(\theta_0)+ \frac{\mathbf{L}}{2}  \,  \hat{S}^2 \,\sum_{k=0}^{K-1}\alpha_k^2+ \frac{\hat{S}^2}{2 M^2} A_K.
%+8L_{\cL}^{(\delta)}\,L\, S_2\big(K_{\mathrm{grad}}^{(\delta)}\big)^2\Sigma_K
%+\frac{8}{M^2}S_1^2\big(K_{\mathrm{grad}}^{(\delta)}\big)^2A_K.
\end{align}
Divide by $A_K$ to obtain the statement.
\end{proof}
%Note that from our  proof of Theorem \ref{lem:Lipschitz_gradient} it also follows that $\theta\mapsto \cL_{L}(\theta)$ is L-smooth. Hence with $\theta_{t+1} =\theta_{t} -\alpha_{t} g_L(\theta_{t})$ we obtain:
%\begin{eqnarray*}
%    \cL_{L}(\theta_{t+1})&\le& \cL_{L}(\theta_{t}) - \alpha_{t}
%\big\langle \nabla\cL_{L}(\theta_{t}),g_L(\theta_{t})\big\rangle
%+\frac{L}{2}\alpha_{t}^2 \|g_L(\theta_{t})\|_2^2.
%\end{eqnarray*}
%Taking the conditional expectation $\EE_G[\cdot |\theta_{t}]$ on both sides yield:
%\begin{eqnarray*}
%    \EE_G[ \cL_{L}(\theta_{t+1})|\theta_{t}] &\le& \cL_{L}(\theta_{t}) - \alpha_{t}
%\big\langle \nabla\cL_{L}(\theta_{t}), \EE_G[g_L(\theta_{t})|\theta_{t}]\big\rangle
%+\frac{L}{2}\alpha_{t}^2 \EE_G[\|g_L(\theta_{t})\|_2^2|\theta_{t}]\\
%&\le& \cL_{L}(\theta_{t}) - \alpha_{t} \| \nabla\cL_{L}(\theta_{t})\|_2^2 + \| \nabla\cL_{L}(\theta_{t})\| \varepsilon_{t}\alpha_{t} 
%+\frac{L}{2}\alpha_{t}^2B^2\\
%&\le & \cL_{L}(\theta_{t}) - \alpha_{t} \| \nabla\cL_{L}(\theta_{t})\|_2^2 + K\varepsilon_{t}\alpha_{t} 
%+\frac{L}{2}\alpha_{t}^2B^2
%\end{eqnarray*}
%Taking the expectation $\EE_G[\cdot]$ on both sides and rearranging terms yields

%\begin{eqnarray*}
%    \EE_G[ \cL_{L}(\theta_{t+1})] - \EE_G[\cL_{L}(\theta_{t})] &\le&   - \alpha_{t} \EE_G[\| \nabla\cL_{L}(\theta_{t})\|_2^2] + K\varepsilon_{t}\alpha_{t} 
%+\frac{L}{2}\alpha_{t}^2B^2
%\end{eqnarray*}
%Summing the inequality over $t=0,\ldots ,T-1$ yields
%\begin{eqnarray*}
%    \EE_G[ \cL_{L}(\theta_{T})] - \EE_G[\cL_{L}(\theta_{0})] &\le&   - \sum_{t=0}^{T-1} \alpha_{t} \EE_G[\| \nabla\cL_{L}(\theta_{t})\|_2^2] + K\sum_{t=0}^{T-1} \varepsilon_{t}\alpha_{t} +\frac{L}{2} B^2\sum_{t=0}^{T-1}\alpha_{t}^2.
%\end{eqnarray*}
%Rearranging and noting that $ \inf_\theta \cL_{L}(\theta) \ge 0 $ yields 
%\begin{eqnarray*}
%  \sum_{t=0}^{T-1} \alpha_{t} \EE_G[\| \nabla\cL_{L}(\theta_{t})\|_2^2] &\le &  \EE_G[\cL_{L}(\theta_{0})]    + K\sum_{t=0}^{T-1} \varepsilon_{t}\alpha_{t} +\frac{L}{2} B^2\sum_{t=0}^{T-1}\alpha_{t}^2.
%\end{eqnarray*}
%Rearranging we obtain
%\begin{eqnarray*}
%    \alpha_k \| \nabla_\theta \hat\cL_{L}(\theta_k)\|^2  &\le & \hat\cL_{L}(\theta_k)- \hat\cL_{L}(\theta_{k+1})+\alpha_k^2 \frac L2 \| \nabla_\theta \hat\cL_{L}(\theta_k)\|^2 
%\end{eqnarray*}
%The L-smoothness of  $\theta\mapsto \hat\cL_{L}(\theta)$  also implies that $\| \nabla_\theta \hat\cL_{L}(\theta_k)\|^2 $ is bounded on $G$ by $C$. Summing the inequality over $k=0,\ldots ,K-1$ yields
%\begin{eqnarray*}
%    \sum_{k=0}^{K-1}\alpha_k \| \nabla_\theta \hat\cL_{L}(\theta_k)\|^2  &\le & \hat\cL_{L}(\theta_0)- \hat\cL_{L}(\theta_{K})+ C \frac L2  \sum_{k=0}^{K-1}\alpha_k^2.
%\end{eqnarray*}
%Letting $T\to\infty$, the right hand side converges against a finite value due to our assumptions. Thus, because of our assumptions on $(\alpha_t)$ we must have $\| \nabla_\theta \cL_{L}(\theta_t)\|\to0$ a.s.\ for $t\to\infty.$

\section{Applications}\label{sec:appl}
We implemented directly the Algorithm \ref{alg:theoretical} with the following set of parameters for discretization of the Fourier domain with a uniform grid with nodes $ u_\ell \;=\; -K + (\ell-1)\,\Delta u$ with $
\Delta u \;=\; \frac{2K}{L}$, for $ \ell=1,\dots,L$. For the smoothing weight we choose Gaussian $w(u)\;=\;e^{-\alpha u^{2}}$ and $\alpha$ is a smoothing parameter (usually around $0.05$ in our examples). 

The policy $(\pi_t)$ is implemented as a feed-forward neural network $f(\theta,s_t,R_t,z_t,t) \sim \pi_t(\cdot|s_t,R_t)$
receiving the four inputs $(s,R,z,t)\in\mathbb{R}^4$.
The inputs are first concatenated and passed through a linear layer of width
$P=256$, followed by layer normalization and a ReLU activation.
This is followed by $K=4$ residual blocks, each consisting of a linear map
$\mathbb{R}^d\!\to\!\mathbb{R}^d$ with layer normalization and ReLU, and a skip
connection.  A final linear layer maps the resulting hidden state to a scalar
output.  The output layer is initialized to zero so that
the initial action distribution is unbiased. No activation function is applied to the final layer. 

For all the examples we run the code until we obtain a loss of order $10^{-3}$, since the algorithm is stochastic and the problem nonconvex, the exact number of steps can vary a lot from run to run, this is especially true for the first example, while the others because of their structure are faster and converge in all runs. Whenever the loss is not below the threshold for $1000$ iterations we rerun the algorithm.   

The exact code for all numerical examples can be found \href{https://github.com/ThanosVasileiadis/Markov-Decision-Processes-of-the-Third-Kind-Learning-Distributions-by-Policy-Gradient-Descent}{\emph{here}} \footnote{https://github.com/ThanosVasileiadis/Markov-Decision-Processes-of-the-Third-Kind-Learning-Distributions-by-Policy-Gradient-Descent} with the exact parameter values. To run our examples we use the HAICORE server, courtesy of the Helmholtz Association. Smaller values of the parameters could run on an M1 Macbook air with 16GB of RAM. 

\subsection{Linear Quadratic Control Problem}
We consider a discrete--time stochastic control problem over a finite horizon \(T\in\mathbb{N}\).
The scalar state \(s_t\) evolves under real--valued actions \(a_t\) from a compact set according to
\begin{equation}\label{eq:dynamics}
  s_{t+1} \;=\; s_t + a_t + \sigma_\varepsilon\,\varepsilon_{t+1}, 
  \qquad t=0,1,\dots,T-1,
\end{equation}
where \(\sigma_\varepsilon=0.1\) is fixed and \((\varepsilon_t)\) are i.i.d.\  standard normal random variables, independent of the initial condition \(s_0=0\).  At each step, a one-stage reward (negative cost) is incurred,
\begin{equation}\label{eq:stage-cost}
  r(s_t,a_t) \;=\; -\tfrac12\bigl(s_t^2 + a_t^2\bigr),
\end{equation}
and the cumulative reward is updated as
\begin{equation}\label{eq:cumulative-cost}
  R_0 := 0,
  \qquad 
  R_{t+1} := R_t + r(s_t,a_t),
  \qquad 
  R_T = \sum_{t=0}^{T-1} r(s_t,a_t) .
\end{equation}
Note that our Assumptions 1,2 and 3 are satisfied and Assumption 4,5 are satisfied on the good event which is sufficient for our considerations. For example for Assumption \ref{ass:rFLip} (i) we have on the good event $G$, $|s_t^2-\tilde s_t^2|\le |s_t-\tilde s_t| 2 B^s_t$.
We are now interested in the distribution of $R_T$. We want to match it with a given distribution. Thus, we consider the lifted MDP:
At time \(t\), the controller observes the current state \(s_t\), the cumulative reward \(R_t\), and the normalized time--to--go
\(\tau_t := (T-t)/T\) which we use instead of $t$ for the practical implementation.
A  randomized Markov policy is a sequence \(\pi=(\pi_t)_{t=0}^{T-1}\) with stochastic kernels
\(\pi_t(\cdot\,|\,s_t,R_t)\) on \(\mathbb{R}\), parametrized by $f(\theta,s,R,z,\tau)$. The distribution of $z$ is a standard Gaussian. 

\paragraph{Target Distribution}
In order to test the algorithm we choose as a target characteristic function one which can be attained by a deterministic policy. More precisely, we choose the optimal LQ feedback control $a_t= -0.5 s_t$ and generate for the target characteristic function $M=100\times 1024$ trajectories \((R_T^{m})_{m=1}^M\) under the dynamics
\eqref{eq:dynamics}--\eqref{eq:cumulative-cost} for $T=10$ and define its empirical characteristic function as
\begin{equation}\label{eq:empirical-policy}
  \widehat{\phi}^*(u)
  \;=\;
  \frac{1}{M}\sum_{m=1}^{M}
  e^{\,iu R_T^{m}} .
\end{equation}
The control objective is to choose \(\pi\) so that the empirical law of \(R_T\) matches the target distribution in the frequency domain as formulated in \eqref{eq: parametrised Loss}.
%, namely to solve
%\begin{equation}\label{eq:empirical-cf-matching}
%  \min_{\pi}
%  \;\;
%  \sum_{k=1}^{N_\omega}
%  K(\omega_k)\,
%  \bigl|
%    \widehat{\phi}_{J_T^\pi}(\omega_k)
%    - \widehat{\phi}_{\star}(\omega_k)
%  \bigr|^2\,\Delta\omega ,
%\end{equation}
%The optimization  is performed over all admissible policy sequences \(\{\pi_t\}_{t=0}^{T-1}\) that are measurable with respect to \((x_t,J_t,\tau_t)\).
%\paragraph{Implementation Details}

\paragraph{Numerical Results and Discussion}
%To obtain a sensible reference distribution of the cumulative cost, we simulate the controlled system \eqref{eq:dynamicsx}-\eqref{eq:cumulative-cost} over a fixed horizon $T=10$ with the same volatility parameter $\sigma_\varepsilon=0.1$.

%By applying the prescribed policy, for instance the optimal LQ feedback control $a_t= -0.5 x_t$, we generate an ensemble of trajectories $\{J_T^{(m)}\}_{m=1}^M$, with $M=100\times 1024$, the resulting empirical distribution serves as the \emph{target} cost distribution for problem \eqref{eq:empirical-cf-matching}. 

The policy gradient algorithm is capable of matching the distribution of cumulative results to the target. The \emph{loss is 0.000384} %and we measure also the Wasserstein distance, $W_1= 0.081414$. To validate further our results we performed a \emph{Kolmogorov-Smirnov} test with \emph{statistic 0.01037 and p-value 0.2803} and a \emph{Von Mises} with \emph{statistic 0.148653 and p-value 0.3937}. 
In Figure \ref{fig:ex1histogram} (a) we see the discrete density of the target versus the learned distribution. In (b) we compare the real and imaginary part of the target and learned characteristic function. 
   
As already explained there is an infinity of optimal action distribution and different runs of the algorithm result in different optimal action distributions as demonstrated by the Figure \ref{fig:different_action_distr}. In this figure we plot the action distribution from our parametrization for different state values $s$ and fixed $R$ and $\tau.$ In most cases the action densities are peaked which means they are close to deterministic.

We emphasise that in accordance with our main Theorem \ref{thm:conv-discrete} our policy gradient algorithm converges to a stationary point and thus it can get stuck in local minima, the results reported above concern only cases where we indeed successfully minimize the loss up to a threshold. There are cases where the algorithm doesn't achieve the minimum loss threshold and gets stuck. We will come back to this point in future work to "unstuck" the algorithm from local minima. 

Finally, the number of trajectories $M=100\times 1024$, seems sufficient to eliminate the bias from the estimator and stabilize the descent, even if it converges to a local minimum.  

\begin{figure}[H]
    \centering
    \begin{subfigure}[b]{0.55\textwidth}
        \includegraphics[width=\textwidth]{cost_distribution}
        \caption{Histogram}
        \label{fig:return_dist2}
    \end{subfigure}
    ~ %add desired spacing between images, e. g. ~, \quad, \qquad, \hfill etc. 
      %(or a blank line to force the subfigure onto a new line)
    \begin{subfigure}[b]{0.4\textwidth}
        \includegraphics[width=\textwidth]{cf_plot}
        \caption{Characteristic Function }
        \label{fig:return_cf2}
    \end{subfigure}
    \caption{Distribution of Rewards}
     \label{fig:ex1histogram}
\end{figure}

\begin{figure}[H]
    \centering
    \begin{subfigure}[b]{0.45\textwidth}
        \includegraphics[width=\textwidth]{action_distributions-7}
        \caption{An optimal action distribution}
        \label{fig:return_dist2-2}
    \end{subfigure}
     %add desired spacing between images, e. g. ~, \quad, \qquad, \hfill etc. 
      %(or a blank line to force the subfigure onto a new line)
    \begin{subfigure}[b]{0.45\textwidth}
        \includegraphics[width=\textwidth]{action_distributions-8}
        \caption{Another optimal action distribution}
    \end{subfigure}
    \caption{Two different optimal action distributions}
    \label{fig:different_action_distr}
\end{figure}

 
\subsection{Targeting Wealth Distributions by Investment}
\label{subse:wealth_example}

We consider a simple discretized Black Scholes model with length of time interval $\Delta t$, i.e. the bond price evolution is given by $B_t=e^{r t \Delta t}$ for $r>0$ and $t\in\NN_0$ and the stock price is given by $x_0=1$ and 
$$ x_{t+1}=x_t \cdot \exp\Big((\mu-\frac12 \sigma^2)\Delta t+\sigma \sqrt{\Delta t} \;\varepsilon_{t+1}\Big)$$
for $\mu\in \mathbb{R}$ and $\sigma>0$ and i.i.d.\ standard normal $(\varepsilon_t)$. If $a_t$ is the amount of money invested in the stock at time $t$ and it is decomposed as $a_t=\pi_t \;s_t$, where $\pi_t\in[0,1]$ is the percentage of our wealth $s_t$ invested in the stock. Then the wealth $s_{t+1}$ at time $t+1$  is given as follows:
$$ s_{t+1}= F(s_t,a_t,\varepsilon_{t+1})  = e^{r  \Delta t} \big( s_t+ a_t y_{t+1}\big)$$
with $$y_{t+1} = e^{-r  \Delta t} \exp\Big((\mu-\frac12 \sigma^2)\Delta t+\sigma \sqrt{\Delta t} \;\varepsilon_{t+1}\Big)-1, \qquad \varepsilon_{t+1} \sim N(0,1),$$
being the excess return over the bond. 
%Let
%\[
%m := (\mu - \tfrac12 \sigma^2 - r)\,\Delta t,
%\qquad
%s := \sigma \sqrt{\Delta t}.
%\]
%Then the variable
%$y_t + 1 = \exp(m + s\,\varepsilon_t)$
%is log-normally distributed with log--parameters $(m,s^2)$. 
As in the motivating example we consider the reward to be non-stationary and set $r_t(s,a)=0$ for $t=0,\ldots,T$ and $r_{T+1}(s,a)=s.$ Thus we are interested in the distribution of the terminal wealth $s_T$ which  is given by 
\begin{equation}
	\label{eq:terminal_wealth}
s_T
= e^{r T \Delta t}
  \Biggl(
        s_0 
      + \sum_{t=1}^{T} 
            a_t \, y_t 
                \exp\!\bigl(-r (t-1)\Delta t \bigr)
  \Biggr).
\end{equation}
Note that when we choose actions from a compact set (which is not really a restriction for applied problems) then Assumptions 1-5 are satisfied on the good event $G$. 


\paragraph{Target Distribution}

For our target distribution, as in the LQ Example, we generate terminal-wealth samples from \eqref{eq:terminal_wealth} under fixed, known investment policies.
We distinguish the following two benchmark scenarios.

\begin{enumerate}
	
	\item \textbf{100\% investment in the stock ($\pi_t=1$).}  
	Here $a_t=s_t$ for all $t$, so the entire wealth is invested in the risky asset at every time step.
	The wealth process follows a multiplicative recursion and can be written as
	\[
	s_T
	= s_0 \prod_{t=1}^T \exp\!\bigl((\mu-\tfrac12\sigma^2)\Delta t
	+ \sigma \sqrt{\Delta t}\,\varepsilon_t \bigr).
	\]
	Consequently, $\log s_T$ is normally distributed and the terminal wealth $s_T$ follows a log-normal distribution,
	\[
	s_T \sim \mathrm{LN}\!\Bigl(
	\log s_0 + T(\mu-\tfrac12\sigma^2)\Delta t,
	\; T\sigma^2\Delta t
	\Bigr).
	\]

	\item \textbf{Random investment fraction ($\pi_t\sim U(0,1)$).}  
	In this case, the investment proportions $\pi_t$ are sampled independently from the uniform distribution on $[0,1]$ and are independent of the market noise $\varepsilon_t$.
	
	The induced distribution of $s_T$ has no closed-form expression; however, we can approximate	this distribution  numerically via Monte--Carlo simulation. 
\end{enumerate}
  
\begin{remark}
\label{re:terminal_wealth}
    Because the terminal wealth depends linearly on the control and multiplicatively on market noise, so randomizing the investment fraction adds negligible variance compared to market fluctuations, making a uniform policy and its mean indistinguishable at the level of the terminal distribution.
\end{remark}

%We set as a target distribution a Gaussian done. To justify this, we recall that Lindberg's condition is satisfied in our example, see \cite[Theorem 27.2]{Billingsley1995}, given a small discretization step, independent noises and smooth bounded controls (which in is also one of the main ingredients of our Theorem \ref{thm:conv-discrete}, at least in the controlled noise regime). 
 



We discretize the horizon $T=1$ into $h=20$ rebalancing dates, hence $\Delta t = T/h$=0.05.
For the Black--Scholes dynamics we use initial wealth $s_0=100$, risk-free rate $r=0.02$, stock drift $\mu=0.06$ and volatility $\sigma=0.40$.
In each epoch, we simulate a Monte--Carlo batch of $M=10^5$ independent trajectories by sampling i.i.d.\ standard Gaussians
$\varepsilon_{t+1}\sim\mathcal N(0,1)$ and iterating the recursion for wealth under the current policy.

%To match terminal-wealth distributions, we work in the Fourier domain on a uniform frequency grid $(u_\ell)_{\ell=1}^{L}$ with $L=301$ over $[-K,K]$ where $K=10.5$, and the characteristic function of terminal wealth from a separate simulation. We compare the resulting empirical characteristic function to the target via a smoothed discrepancy with kernel bandwidth $\alpha=0.025$. We perform  $k=201$ iterations of the descent using a fixed learning rate $\texttt{lr}=10^{-4}$, with gradients obtained by automatic differentiation through the simulation.

\paragraph{Numerical Results and Discussion} 

For the first  case we see in Figure \ref{fig:wealth_dist_case_2} and \ref{fig:cf_dist_wealth_case_2}  the histogram of the wealth distributions (target and learned) as well as the real and imaginary part of the target and learned characteristic functions. Both plots show good fits.  From Figure  %\ref{fig:optimal_action_dist_wealth_case_1} and 
\ref{fig:optimal_action_dist_wealth_case_2} it is clear that the optimal action distribution produced by the algorithm is clearly deterministic and close to the theoretical one. 

For the second case, we plot first in Figure \ref{fig:wealth_dist_case_3} the comparison of the target density (blue) and the density of the terminal wealth learned by the algorithm (red). In Figure \ref{fig:cf_dist_wealth_case_3} we compare the real and imaginary part of the target and learned characteristic functions. Both plots show a good fit.
The optimal action distribution shown in Figure \ref{fig:optimal_action_dist_wealth_case_3} is as expected by Remark \ref{re:terminal_wealth} close to the mean of the uniform $(0,1)$ and almost deterministic, i.e. with a very small variance. This discrepancy could be due to precision error of the calculation. We investigate further the accuracy of the approximation in the next example where we have an exact analytically traceable unique distribution for the optimal controls.    

\textbf{Case 1) $\pi_t=1$}\\



\begin{figure}[H]
    \centering
    \begin{subfigure}[b]{0.45\textwidth}
        \includegraphics[width=\textwidth]{terminal_wealth_distribution-6}
        \caption{Histogram}
        \label{fig:wealth_dist_case_2}
    \end{subfigure}
    ~ %add desired spacing between images, e. g. ~, \quad, \qquad, \hfill etc. 
      %(or a blank line to force the subfigure onto a new line)
    \begin{subfigure}[b]{0.4\textwidth}
        \includegraphics[width=\textwidth]{cf_comparison-9}
        \caption{Characteristic Functions}
        \label{fig:cf_dist_wealth_case_2}
    \end{subfigure}
    \caption{Distribution of terminal wealth}
    \end{figure}

\begin{figure}[H]
	\centering
	\includegraphics[width=0.5\textwidth]{action_distributions-9.png}
	\caption{Optimal actions distributions for Different wealth states}
	\label{fig:optimal_action_dist_wealth_case_2}
\end{figure}

\textbf{Case 2) $\pi_t\sim U(0,1) $}\\

\begin{figure}[H]
    \centering
    \begin{subfigure}[b]{0.45\textwidth}
        \includegraphics[width=\textwidth]{terminal_wealth_distribution-2}
        \caption{Histogram}
        \label{fig:wealth_dist_case_3}
    \end{subfigure}
    ~ %add desired spacing between images, e. g. ~, \quad, \qquad, \hfill etc. 
      %(or a blank line to force the subfigure onto a new line)
    \begin{subfigure}[b]{0.4\textwidth}
        \includegraphics[width=\textwidth]{cf_comparison-7}
        \caption{Characteristic Functions}
        \label{fig:cf_dist_wealth_case_3}
    \end{subfigure}
    \caption{Distribution of terminal wealth}
    \end{figure}


\begin{figure}[H]
	\centering
	\includegraphics[width=0.5\textwidth]{action_distributions-4.png}
	\caption{Optimal actions distributions for Different wealth states}
	\label{fig:optimal_action_dist_wealth_case_3}
\end{figure}

\subsection{An Example with Compact Support}
This example is related to the first one. Our goal is to exemplify a rather unique phenomenon. Due to the structure of the cost function we can identify a unique control distribution that stirs the discounted rewards towards the target. In contrast with Example \ref{subse:wealth_example}, here the control enters inside a nonlinear function and together with the noise on the dynamics shape the reward distribution.

The policy gradient can indeed recover the unique control distribution when the length of the action interval permits unique identifiability of policies or otherwise only by enforcing the symmetry of the control distribution a posteriori. 

In concrete, we consider linear dynamics with $T=1$, a randomized (relaxed) control $a_0$ and additive Gaussian noise:
\begin{equation}
\label{eq:dynamics2}
 s_1 = s_0 + a_0 + \varepsilon_1, \qquad \varepsilon_1 \sim \mathcal N(0,\sigma_\varepsilon^2),
\end{equation}
where $s_0\in\mathbb R$ is known and $a_0$ is a control random variable, independent of $\varepsilon_1$.
Define the bounded performance variables $r_0\equiv 0$ and
\begin{equation}
\label{eq:r-def}
 r_1(s_1,a_1) := \cos(s_1) \in [-1,1],
\end{equation}
thus, we are interested in the law of $R_2=\cos(s_0+a_0+\varepsilon_1).$ Obviously we can restrict $a_0$ to a compact set $I$ and our Assumption 1-5 are satisfied.
\paragraph{Target Distribution}
Our \emph{target} distribution is the Epanechnikov law on $[-1,1]$ with density
\begin{equation}
\label{eq:epa-density}
 f_E(x) = \frac{3}{4}\Bigl(1 - x^2\Bigr)\,1_{\{|x|\le 1\}},
\end{equation}
its characteristic function admits the closed form
\begin{equation}
\label{eq:phiE}
 \phi_E(u) = \frac{3}{u^3}\Bigl(\sin(u ) - u \cos(u)\Bigr), \qquad \phi_E(0)=1.
\end{equation}

We seek a distribution of the randomized control $a_0$ such that the distribution of $R_2$ matches the Epanechnikov law. 
%\paragraph{Compact action space.} 
We assume $a_0$ takes values in a compact interval $I$ of length $|I|<\pi$, so that its law is uniquely determined by its Fourier coefficients on $\mathbb Z$. This ensures that the Fourier coefficients $\psi_A(k):=\EE\big[e^{ik a_0}\big]$ with $\,k\in\ZZ $ fully characterize the distribution of $a_0$ and its  characteristic function by 
\begin{equation}
\phi_A(t) = \int_I e^{itx}p_{a_0}(x)\,dx = \sum_{k=-\infty}^{\infty} \psi_A(k)\,\int_I e^{i(t-k)x}\,dx.
\end{equation}

\paragraph{Exact matching condition}
With the help of the Jacobi--Anger identity we can compute the characteristic function of $R_2$ as follows:
%The Jacobi--Anger identity for $z,\vartheta\in\mathbb R$ is
%\begin{equation*}
% e^{iz\cos\vartheta} = \sum_{k=-\infty}^{\infty} i^k J_k(z)e^{ik\vartheta},
%\end{equation*}
%where $J_k$ is the $k$-th Bessel function of the first kind. Applying this with $z=u V $ and $\vartheta=s_1$ gives
\begin{align*}
\phi_{R_2}(u) &= \EE[e^{iu R_2}] = \sum_{k\in\mathbb Z} i^k J_k(u V)\EE[e^{iks_1}] = \sum_{k\in\mathbb Z} i^k J_k(u V )e^{iks_0}e^{-\frac12 k^2\sigma^2}\psi_A(k)
\end{align*}
where $J_k$ is the $k$-th Bessel function of the first kind. Thus the matching equation for the characteristic functions is
\begin{equation}
\label{eq:matching-functional}
 \sum_{k\in\mathbb Z} i^kJ_k(u V)e^{iks_0}e^{-\tfrac12 k^2\sigma^2}\psi_A(k) = \phi_E(u), \quad \forall u\in\mathbb R.
\end{equation}

%When the action space is a compact $I$, the moments $\psi_A(k)$ uniquely determine the full characteristic function
%\begin{equation}
%\phi_A(t) = \int_I e^{itx}p_{a_0}(x)\,dx = \sum_{k=-\infty}^{\infty} %\psi_A(k)\,\int_I e^{i(t-k)x}\,dx,
%\end{equation}
%so the solution of \eqref{eq:matching-functional} provides full identifiability of $a_0$.

\paragraph{Truncation and linear system}
Since $J_k$ and $e^{-k^2\sigma^2/2}$ decay rapidly, we truncate to $|k|\le K$ and select a finite frequency grid $\{u_\ell\}_{\ell=1}^L\subset[-K,K]$. The truncated system is
\begin{equation}
\label{eq:trancated_linear_system}
 J_0(u_\ell V) + 2\sum_{k=1}^{K} J_k(u_\ell V)e^{-\frac12k^2\sigma^2}x_k = \phi_E(u_\ell), \quad x_k=Re(i^ke^{iks_0}\psi_A(k)).
\end{equation}
Note that the symmetry $\psi_A(-k) = \overline{\psi_A(k)}$ and $J_{-k} = (-1)^k J_k$ makes the imaginary parts cancel. Solving for $x=(x_1,\dots,x_K)$ yields a least-squares (or exact solution). In addition both left and right-hand side are even, forcing $x_k=0$ for all odd $k$. Then we reconstruct
\begin{equation}
 \psi_A(k) = i^{-k}e^{-iks_0}x_k, \qquad |\psi_A(k)|\le 1,\quad \psi_A(-k)=\overline{\psi_A(k)}.
\end{equation}

The density on $I$ (of width $<\pi$) follows by Fourier inversion:
\begin{equation}
 p_{a_0}(x) = \frac{1}{2\pi}\Big(1+2\sum_{k=1}^{K}Re(\psi_A(k)e^{-ikx})\Big),\quad x\in I.
\end{equation}

Finally, we notice that $p_{a_0}$ is $\pi$-periodic since all odd Fourier modes $\psi_A(k)$ are zero. 

\paragraph{Numerical Results and Discussion}
First for the numerical solution of \eqref{eq:trancated_linear_system} we fix $K=16$ and $L>>K$ in particular $L=8001$ (for increased accuracy in our policy gradient since we used the same grid), and thus implement a least square solution. If we write the right-hand side  of \eqref{eq:trancated_linear_system} in matrix form we need the corresponding matrix to have full rank - a condition easily verifiable in our case - in order to have a unique solution. For the MC samples we used $M=72 \times 1024$ trajectories per iteration. 

We run two sets of experiments one on the interval $I_1=[0,\pi]$ to have uniquely identifiable actions and one on the interval $I_2=[-\pi,\pi]$ to try and identify uniquely the optimal policy enforcing the symmetry of the distribution. 

\noindent
\textbf{Case $I_1=[0,\pi]$}\\

\begin{figure}[H]
    \centering
    \begin{subfigure}[b]{0.45\textwidth}
        \includegraphics[width=\textwidth]{return_distribution-4_2}
        \caption{Histogram}
        \label{fig:return_dist2-3}
    \end{subfigure}
    ~ %add desired spacing between images, e. g. ~, \quad, \qquad, \hfill etc. 
      %(or a blank line to force the subfigure onto a new line)
    \begin{subfigure}[b]{0.4\textwidth}
        \includegraphics[width=\textwidth]{cf_match-3_2}
        \caption{Characteristic Function }
        \label{fig:return_cf2-2}
    \end{subfigure}
    \caption{Distribution of Rewards}
\end{figure}

\begin{figure}[H]
    \centering
    \begin{subfigure}[b]{0.45\textwidth}
        \includegraphics[width=\textwidth]{actions_vs_solver_hist-4_2}
        \caption{Histogram}
        \label{fig:action_dist2}
    \end{subfigure}
    ~ %add desired spacing between images, e. g. ~, \quad, \qquad, \hfill etc. 
      %(or a blank line to force the subfigure onto a new line)
    \begin{subfigure}[b]{0.4\textwidth}
        \includegraphics[width=\textwidth]{moment_magnitude-4_2}
        \caption{Fourier modes $\psi_A(k)$}
        \label{fig:action_fourier_2}
    \end{subfigure}
    \caption{Action distribution without projecting}
\end{figure}
First note that the density of learned distribution is rather close to the target distribution, cf. Figure \ref{fig:return_dist2-3}. In Figure \ref{fig:return_cf2-2} we compare the real and imaginary part of the learned and target characteristic function. Since we restrict to $I_1=[0,\pi]$, the imaginary part is not zero in the learned case, as it is in the target case.
However, here the policy gradient algorithm can indeed find the unique action distribution. An attentive reader could notice that in Figure \ref{fig:action_fourier_2} the odd modes are not zero as they should be. This is due to the fact that we have a $\pi$-periodic density and we have samples only on the one side of the density.  To reveal the true symmetric density we need to have a look at the interval $I_2=[-\pi,\pi]$ as we do in the next experiment. 
%A very simple justification of this fact can be done by the solving 
%\begin{align*}
%	&\EE[e^{iA}]=0 \\
%	&\EE[\cos(A)+i\sin(A)]= \EE[\mathit{1}_{\{A=a, \, a\in[0,\pi] \}} \cos(A)+i\sin(A)]+\EE[\mathit{1}_{\{A=a-\pi,\, a\in[0,\pi]\}} \cos(A)+i\sin(A)]=0 %    \end{align*}
%where $A$ is the random variable representing the control and when the distribution assigns equal weigh to each branch of period $\pi$. 

So when only one period is used the symmetry cannot be obtain in terms of Fourier modes but the Wasserstein distance is effectively minimal  $ {W_1=0,01302292}$.

 
\noindent
\textbf{Case $I_2=[-\pi,\pi]$}\\
Here the Solver solution for \eqref{eq:trancated_linear_system} gives two periods of the $\pi$-periodic density. The Policy gradient is not capable to identify correctly the complete periodic density since for the reward there is no way to differentiate between periodic actions and stays concentrated on the one of the two branches of the density (since the algorithm is stochastic the outcome may vary from run to run, meaning that sometimes the left lobe is selected and other times the right, cf. Figures \ref{fig:return_dist2-4} and \ref{fig:return_cf2-3}  for the densities and characteristic functions of the rewards and Figures \ref{fig:action_dist2-2} and \ref{fig:action_fourier} for the action densities and Fourier modes of the action distribution.
  

\begin{figure}[H]
    \centering
    \begin{subfigure}[b]{0.45\textwidth}
        \includegraphics[width=\textwidth]{return_distribution-4}
        \caption{Histogram}
        \label{fig:return_dist2-4}
    \end{subfigure}
    ~ %add desired spacing between images, e. g. ~, \quad, \qquad, \hfill etc. 
      %(or a blank line to force the subfigure onto a new line)
    \begin{subfigure}[b]{0.4\textwidth}
        \includegraphics[width=\textwidth]{cf_match-3}
        \caption{Characteristic Function }
        \label{fig:return_cf2-3}
    \end{subfigure}
    \caption{Distribution of Rewards}
    \end{figure}

\begin{figure}[H]
    \centering
    \begin{subfigure}[b]{0.45\textwidth}
        \includegraphics[width=\textwidth]{actions_vs_solver_hist-5_2}
        \caption{Histogram}
        \label{fig:action_dist2-2}
    \end{subfigure}
    ~ %add desired spacing between images, e. g. ~, \quad, \qquad, \hfill etc. 
      %(or a blank line to force the subfigure onto a new line)
    \begin{subfigure}[b]{0.4\textwidth}
        \includegraphics[width=\textwidth]{moment_magnitude_interval}
        \caption{Fourier modes $\psi_A(k)$}
        \label{fig:action_fourier}
    \end{subfigure}
    \caption{Action distribution without projecting}
\end{figure}

Nevertheless, if we split equally and project the samples to the opposite interval after training we can recover the correct periodic density from the inflated one of Figure \ref{fig:action_dist2}, see Figures \ref{fig:gull} and \ref{fig:tiger}.

\begin{figure}[H]
    \centering
    \begin{subfigure}[b]{0.45\textwidth}
        \includegraphics[width=\textwidth]{actions_vs_solver_hist-5}
        \caption{Histogram}
        \label{fig:gull}
    \end{subfigure}
    ~ %add desired spacing between images, e. g. ~, \quad, \qquad, \hfill etc. 
      %(or a blank line to force the subfigure onto a new line)
    \begin{subfigure}[b]{0.4\textwidth}
        \includegraphics[width=\textwidth]{moment_magnitude-4}
        \caption{Fourier modes}
        \label{fig:tiger}
    \end{subfigure}
    \caption{Action distribution with projection}
\end{figure}

\begin{remark}[about the fit of the action distributions]
	To test the best possible fit for the action distribution we augment the loss function with an additional matching objective $\sum_{j=0}^{K=16}\vert \psi_A(k)- \hat\psi_A(k)w(u_\ell)\vert^2 $ for the Fourier modes that come from the numerical solution of \eqref{eq:trancated_linear_system}. As one can see in Figure \ref{fig:Policy_vs_solver} the periodic densities are very close and thus the expressivity of the ANN enough. We don't attempt here an exhaustive optimization of the architectures and the hyper-parameters since the goal is a proof of concept. 
		\begin{figure}[!h]
		\centering 
		\includegraphics[width=0.8\textwidth]{remark}
\caption{Policy Distributions comparison}
		\label{fig:Policy_vs_solver}
	\end{figure}
\end{remark}

\subsection{Writing a Classical MDP as a Special Case Distributional MDP}
\label{se:forzen_lake}
We give another example, the standard \texttt{FrozenLake} MDP from \cite{towers2025gymnasiumstandardinterfacereinforcement}, a pleasant variant of \cite[Example 3.5]{SuttonBarto}. 
The episodic return of the MDP satisfies
\[
R_T \in \{0,1\},
\qquad
R_T = 1 \text{ if the goal is reached,}
\quad
R_T = 0 \text{ otherwise}.
\]
Let
\[
p_\sigma := \mathbb{P}^\sigma(R_T = 1)
\]
denote the success probability under policy $\sigma$. Then
\[
\mathrm{Law}^\sigma(R_T)
    = p_\sigma\,\delta_1 + (1-p_\sigma)\,\delta_0.
\]
Choose the target distribution as the ideal success law,
$\mu_{\mathrm{target}} = \delta_1.$
If we consider the squared $L^2$ distance between characteristic functions,
\[
\cL(\sigma)
:= \int_{\mathbb{R}}
\big|\varphi^\star(u)-\EE^\sigma[e^{iuR_T^\theta}] \big|^2
\, w(u)\, du,
\]
where $w$ is a nonnegative weight function. A direct computation yields
\[
\cL(\sigma)
= (1 - p_\sigma)^2
\int_{\mathbb{R}} |e^{iu} - 1|^2 w(u)\, du
=: C_w (1 - p_\sigma)^2,
\]
with $C_w \in (0,\infty)$ independent of $\sigma$ 


%The TV distance between the policy-induced return distribution and the target is:
%\[
%\begin{aligned}
%\big\|\mathrm{Law}_\sigma(R_T)-\delta_1\big\|_{\mathrm{TV}}
%    &= \frac12\Big( 
%        |p_\sigma-1| + |(1-p_\sigma)-0|
%    \Big)\\[1ex]
%    &= 1 - p_\sigma.
%\end{aligned}
%\]

%Therefore,
%\[
%\min_\sigma 
%    \big\|\mathrm{Law}_\sigma(R_T)-\delta_1\big\|_{\mathrm{TV}}
%    \quad\Longleftrightarrow\quad
%\max_\sigma p_\sigma
%    = \max_\sigma \mathbb{E}_\sigma[R_T].
%\]

%Thus, in the Frozen Lake classical MDP, the TV-based distribution matching objective is  \emph{exactly equivalent} to maximizing the expected return. The distributional formulation  remains intact, while the divergence specializes to the classical control objective.

%\begin{remark}
%	Conceptually our approach can recover the results of classical control but the exact implementation of the algorithm escapes the scope of the article (since the modified loss is no longer differentiable wrt the control and we need to introduce a different formulation for controls having a density). \footnote{We will deal in a follow up work with this case.} 
%\end{remark}

%If we consider the squared $L^2$ distance between characteristic functions,
%\[
%\cL(\sigma)
%:= \int_{\mathbb{R}}
%\big|\varphi_\sigma(\omega) - \varphi^\star(\omega)\big|^2
%\, w(\omega)\, d\omega,
%\]
%where $w$ is a nonnegative weight function. A direct computation yields
%\[
%\cL(\sigma)
%= (1 - p_\sigma)^2
%\int_{\mathbb{R}} |e^{i\omega} - 1|^2 w(\omega)\, d\omega
%=: C_w (1 - p_\sigma)^2,
%\]
%with $C_w \in (0,\infty)$ independent of $\sigma$ and the explicit identity
%\[
%\cL(\sigma) = C_w \, \|\text{Law}_\sigma(R_T) - \delta_1\|_{\mathrm{TV}}^2,
%\]
Thus, minimizing the characteristic function loss is  equivalent to
\[
\min_\sigma \cL(\sigma)
\quad \Longleftrightarrow \quad
\max_\sigma p_\sigma
\quad \Longleftrightarrow \quad
\max_\sigma \mathbb{E}_\sigma[R_T].
\]

Of course, this theoretical framework has to be adapted in practice to design practical algorithms since matching a Dirac is a difficult problem and also the issue of exploration has to be addressed. 

%Before we explain this, let us note the following: Our target function is $\cL(\theta) $ and the derivative of it is according to \eqref{eq:def:Gradient_Loss}
%\begin{align}
%    \nabla_\theta \mathcal{L}(\theta) =
%    2 \int_{\mathbb{R}} \operatorname{Re} \left[ \overline{\left(\EE^\theta[e^{iu R_T}]- \varphi^\star(u)   \right)}\cdot \nabla_\theta \varphi_{\theta}(u) \right] w(u) du,
%\end{align}
%, we obtain that in the algorithm (except for the discretization of the integral) we replace $ \nabla_\theta \mathcal{L}(\theta) $ by its empirical counterpart
%\begin{align}
%    2 \int_{\mathbb{R}} \operatorname{Re} \left[ \overline{\left(\ \psi_N^\theta(u)- \varphi^\star(u)   \right)}\cdot \nabla_\theta  \psi_N^\theta(u)\right] w(u) du.
%\end{align}
%But it is easy to see that this is the same as the derivative of the empirical target function
%\begin{equation}
%	 \int_{\mathbb{R}} \left|    \psi_N^\theta(u)  -\varphi^\star(u) \right|^2 w(u) \, du.
%\end{equation}
%However in this setting the empirical target function has the explicit representation (see e.g. Epps and Pulley (1983), BH88)  :
%\begin{eqnarray*}
%\mathcal{L}(\theta)  &=& \int_{\mathbb{R}} \left|    \psi_N^\theta(\omega)  -\varphi^\star(\omega) \right|^2 w(\omega) \, d\omega\\
%&=& \frac{1}{N^2} \sum_{j,k=1}^N \exp(-\frac12 |R_T^j-R_T^k|^2) -\sqrt{2}  \frac1N  \sum_{j=1}^N \exp(-1/4 |R_T^j|^4) + \frac{1}{\sqrt{3}}.
%\end{eqnarray*} 

\vspace*{1cm}
\noindent
{\bf  Acknowledgment:}

{This work is supported by the Helmholtz Association Initiative and Networking Fund on the HAICORE@KIT partition. 
Apart from this, the authors received no financial support for the research, authorship, and/or publication of this article and have no competing interests to decalre.}
 
%   All acknowledgements go at the end of the paper before appendices and references.Moreover, you are required to declare funding (financial activities supporting thesubmitted work) and competing interests (related financial activities outside the submitted work).More information about this disclosure can be found on the JMLR website.}

\section{Appendix}
In what follows we use as norm the $L^1$-norm i.e. for $x\in\RR^n$ we have $\|x\|=|x_1|+\ldots +|x_n|.$ The $L^2$ norm is denoted by $\|x\|_2$. Note that we always  have $\|x\|_2 \le \|x\|\le \sqrt{n}\|x\|_2. $
\subsection{The Gaussian Concentration Inequality for Lipschitz Functions}\label{sec:concentration}
 The next result follows from the Gaussian Isoperimetric Inequality and will be used to control the fluctuations of our network's output, for details see \cite{ledoux1994isoperimetry}.
\begin{lemma}[Gaussian Concentration for Lipschitz Functions] 
\label{lem:gaussian-concentration}
Let $Z\in\mathbb{R}^d$ be a Gaussian random vector $Z\sim \mathcal{N}(0,\sigma_z^2 I_d)$ with standard deviation $\sigma_z$ in each direction. Suppose $f: \mathbb{R}^d \to \mathbb{R}$ is a Lipschitz continuous  function with constant $L$. Then for any $\epsilon > 0$, the deviation of $f(Z)$ from its mean is exponentially small
\[
\PP\Big(|f(Z) - \mathbb{E}[f(Z)]| > \epsilon\Big) ~\le~ 2 \exp\!\Big(-\frac{\epsilon^2}{2\,L^2\,\sigma_z^2}\Big)\,. 
\] 
In particular, with probability at least $1-\delta$, one has 
\[
|f(Z)-\mathbb{E}[f(Z)]| \le L\,\sigma_z \sqrt{2\ln(2/\delta)}\,.
\] 
\end{lemma}

\subsection{Bounds}\label{sec:bounds}
In what follows we need bounds on several quantities. We consider everything on $\Theta\times G.$
From \eqref{eq:s_bound}  we obviously have constants $B^s_t$ and $B^z$ such that for all $t:$
\begin{equation}
    |s_t^\theta| \le B^s_t,\quad  |z_t| \le B^z.
\end{equation}
We set  $B^s = \max_{t=0,1,\ldots,T} B_t^s.$
Since $\Theta$ is bounded we can define 
\begin{equation}
    B^\theta = \max_{\theta\in\Theta} \|\theta\|
\end{equation}
Now we can prove by induction:

\begin{lemma}\label{lem:boundssa}
    There exist constants $B_t^R,B_t^a>0, t=0,\ldots,T$ such that
    \begin{equation}
         |R_t^\theta(\omega)| \le B_t^R,\quad  |a_t^\theta(\omega)|\le B_t^a, \quad (\theta,\omega)\in\Theta\times G
    \end{equation}
    and we define $B^R := \max_{t=0,1,\ldots,T} B_t^R, B^a := \max_{t=0,1,\ldots,T} B_t^a.$ 
\end{lemma}

\begin{proof}
    The proof is by induction over $t$. For $t=0$ the statement is obvious. Suppose the bounds are correct up to time $t-1.$ We obtain by the continuity of $r$
   \begin{equation}
    |R_t^\theta| \le \sum_{k=0}^{t-1} |r(s_k^\theta,a_k^\theta)| \le \sum_{k=0}^{t-1} \max_{|s|\le B^s_k, |a|\le B^a_{k}}|r(s,a)| =: B^R_t. 
\end{equation} 

Further we have by continuity of $f$
\begin{equation}
    |a_t^\theta| = |f(\theta,s_t^\theta,R_t^\theta,z_t,t)|\le 
    \max_{\theta\in\Theta, |s|\le B_t^s, |R|\le B_t^R, |z|\le B^z, k\le T} |f(\theta,s,R,z,k)| =: B^a_t
\end{equation}
which proves the statement.
\end{proof}
In particular with Assumption \ref{as:1} and on $\Theta\times G$ we can define 
$$\sigma_{max} :=  \max_{\|x\|\le B^s+B^R+B^z+T,\theta\in\Theta}\sigma(W^1 x+b^1),\quad \sigma'_{max} := \sup_{y\in\RR} \sigma'(y), \quad \sigma''_{max} := \sup_{y\in\RR} \sigma''(y)$$
and $B^x := B^s+B^R+B^z+T.$
The proof of the next lemma follows directly from our assumptions.

\begin{lemma}
    Under  Assumption \ref{ass:rFLipgrad}  and with Lemma \ref{lem:boundssa} we obtain constants $B_F^i, B_r^i,B_\nabla^{f,\theta},B_\nabla^{f,i}  $ such that on $\Theta\times G$ for all $t$:
    \begin{equation}
    |\partial_i F(s^\theta_t(\omega),a^\theta_t(\omega),\varepsilon(\omega))| \le  B_F^i, \quad  |\partial_i r(s^\theta_t(\omega),a^\theta_t(\omega))| \le  B_r^i,\quad i=s,a
\end{equation}
\begin{equation}
    \|f_\theta(\theta,x^\theta_t(\omega))\| \le B_\nabla^{f,\theta},\quad |\partial_i f(\theta,x^\theta_t(\omega))| \le  B_\nabla^{f,i}, \quad i=s,R.
\end{equation}
\end{lemma} 
%\begin{equation}
%    |\partial_i F(s,a,\varepsilon)| \le \max_{|s|\le B^s, |a|\le B^a, |\varepsilon|\le \Gamma}|\partial_i F(s,a,\varepsilon)| =: B_F^i, \quad  |\partial_i r(s,a)| \le \max_{|s|\le B^s, |a|\le B^a}|\partial_i r(s,a)| =: B_r^i.
%\end{equation}
%Further from Lemma \ref{lem:NNLipschitz} below it follows that
%\begin{equation}
%    |f_\theta(\theta,x)| \le B_\nabla^{f,\theta},\quad |f_i(\theta,x)| \le \max_{|\theta'|\le B^\theta, |s'|\le B^s, |R'|\le B^R, |z'|\le B^z}|f_i(\theta',x')| =: B_\nabla^{f,i}.
%\end{equation}
Last but not least we need bounds on $\nabla_\theta s_t^\theta, \nabla_\theta a_t^\theta, \nabla R_t^\theta. $ 
\begin{lemma}\label{lem:bound1}
    Under  Assumption \ref{ass:rFLipgrad} we obtain constants $B_\nabla^{s}, B_\nabla^a, B_\nabla^R$ such that on $\Theta\times G$ for all $t$:
    \begin{equation}
        \|\nabla_\theta s_t^\theta(\omega)\| \le B_\nabla^{s},\;  \|\nabla_\theta a_t^\theta(\omega)\| \le B_\nabla^a,\;  \|\nabla_\theta R_t^\theta(\omega)\| \le B_\nabla^R
    \end{equation}
\end{lemma} 

\begin{proof}
We have that 
\begin{eqnarray}\label{eq:nablas0}
\nabla_\theta s_t^\theta &=& \partial_sF(s_{t-1}^\theta,a_{t-1}^\theta,\varepsilon_t) \nabla_\theta s_{t-1}^\theta + \partial_aF(s_{t-1}^\theta,a_{t-1}^\theta,\varepsilon_t) \nabla_\theta a_{t-1}^\theta \\\label{eq:nablaa}
\nabla_\theta a_t^\theta &=& f_\theta(\theta,s_{t}^\theta,R_{t}^\theta,z_t,t)+ \partial_s f(\theta,s_{t}^\theta,R_{t}^\theta,z_t,t) \nabla_\theta s_{t}^\theta + \partial_R f(\theta,s_{t}^\theta,R_{t}^\theta,z_t,t) \nabla_\theta R_{t}^\theta\\
\nabla_\theta R_t^\theta &=& \sum_{k=0}^{t-1}  \partial_sr(s_k^\theta,a_k^\theta) \nabla_\theta s_k^\theta + \partial_ar(s_k^{\theta},a_k^{\theta}) \nabla_\theta a_k^\theta.
\end{eqnarray}
From this it follows inductively that $\|\nabla_\theta s_t^\theta\|\le B_\nabla^{s,t}, \|\nabla_\theta a_t^\theta\| \le B_\nabla^{a,t}, \|\nabla R_t^\theta\|\le B_\nabla^{R,t}$ and the constants satisfy the following recursion:
\begin{eqnarray*}
    B_\nabla^{s,t} &=& B_F^s B_\nabla^{s,t-1}+ B_F^a B_\nabla^{a,t-1}\\
    B_\nabla^{R,t} &=& \sum_{k=0}^{t-1} B_r^s B_\nabla^{s,k}+ B_r^a B_\nabla^{a,k}\\
    B_\nabla^{a,t} &=& B_\nabla^{f,\theta}+ B_\nabla^{f,s} B_\nabla^{s,t} + B_\nabla^{f,R}B_\nabla^{R,t}.
\end{eqnarray*}
As before we define $B_\nabla^{s}=\max_t B_\nabla^{s,t}, B_\nabla^{a} =\max_t B_\nabla^{a,t}, B_\nabla^{R}=\max_t B_\nabla^{R,t}$
\end{proof}

\subsection{Lipschitz properties}
The proof of the following result is obvious, see e.g. \cite{hinderer2005lipschitz}, Lemma 2.1. 
\begin{lemma}\label{lem:lipschitzproperties}
Let $M,M'$ and $M''$ be metric spaces.
\begin{itemize}
\item[a)] If $f,g:M\to M'$ are Lipschitz continuous with constant $L_f$ and $L_g$ respectively, then  $f +g$ is Lipschitz continuous with constant  $L_f+L_g$.
\item[b)] If $f,g:M\to \RR$ are Lipschitz continuous with constant $L_f$ and $L_g$ respectively and  $| f|\le B_f$ and $| g|\le B_g$ on $M$ respectively, then $f \cdot g$ is Lipschitz continuous with constant $B_g L_f+B_f L_g. $
\item[c)] Let $f:M\times M'\to M''$. If $f$ is  Lipschitz continuous separately in the components, i.e.
$$ \|f(x,y)-f(\tilde x,y)\| \le L_1 \|x-\tilde x\|, \quad \|f(x,y)-f(x,\tilde y)\| \le L_2 \|y-\tilde y\|,$$
for suitable $L_1,L_2,$ then for $z=(x,y), \tilde z=(\tilde x,\tilde y)$ we have with $L := \max\{L_1,L_2\}.$
$$ \|f(z)-f(\tilde z)\| \le L \|z-\tilde z\|$$
\item[d)] If $f:M\to M'$ and $g:M'\to M''$ are Lipschitz continuous with constant $L_f$ and $L_g$ respectively, then $g\circ f$ is again Lipschitz continuous with constant $L_f\cdot  L_g$.
\end{itemize}
\end{lemma}

Next we investigate the Lipschitz properties of our neural net. 
\begin{lemma}\label{lem:NNLipschitz}
Let $f : \RR^{11}\to \RR$ be the neural network defined in \eqref{eq:definition_NN_formula}, i.e.
$$ f(\theta,x) = w^2 \sigma(W^1 x+b^1)+b^2$$
for $\theta=(W^1,w^2,b^1,b^2)\in \RR^7$ and $x=(s,R,z,t)\in\RR^4.$ 
On $\Theta\times G$ we obtain
\begin{itemize}
\item[a)] $f$ is jointly  Lipschitz-continuous in $(\theta,x)$ i.e.
$$|f(\theta,x)-f(\tilde\theta,\tilde x)| \le L_f \|(\theta,x)-(\tilde\theta,\tilde x)\| \mbox{ for all } \omega \in G \mbox{ and } \theta,\tilde\theta\in\Theta$$ with constant
$L_f := \max\{ 1, \sigma_{max}, B^{\theta}\sigma'_{max} ,B^{\theta} B^{x} \sigma'_{max}, (B^\theta)^2 \sigma'_{max}\}. $
\item[b)] $\nabla_\theta f$ is jointly  Lipschitz-continuous in $(\theta,x)$  i.e.
$$|\nabla_\theta f(\theta,x)-\nabla_\theta f(\tilde\theta,\tilde x)| \le L_f^\nabla \|(\theta,x)-(\tilde\theta,\tilde x)\| \mbox{ for all } \omega \in G \mbox{ and } \theta,\tilde\theta\in\Theta$$
with constant $L_f^\nabla = L_{b^1}^\nabla+L_{w^2}^\nabla+L_{W^1}^\nabla$
\begin{eqnarray*}
    L_{b_1}^\nabla & :=& \max\{\sigma'_{max},  B^{\theta}\sigma''_{max},  (B^\theta)^2\sigma'_{max}, B^{\theta} B^{x} \sigma''_{max}\}\\
     L_{w^2}^\nabla & :=& \sigma'_{max} \max\{ 1,B^x, B^{\theta}\}\\
 L_{W^1}^\nabla & :=& \max\{\sigma'_{max} B^\theta,  B^x  \sigma'_{max}, B^{\theta} (B^{x})^2\sigma''_{max},  B^{\theta} B^{x} \sigma''_{max},  (B^{\theta})^2 B^{x} \sigma''_{max}\}.
\end{eqnarray*}
\end{itemize}
\end{lemma}

\begin{proof}
a) According to Lemma \ref{lem:lipschitzproperties} it is enough to show Lipschitz continuity separately in the components. We obtain
\begin{eqnarray*}
|f(\theta,x)-f(\tilde\theta,x)| &\le& |b^2-\tilde b^2| + |w^2-\tilde w^2| \sigma(W^1x+b^1) + |\tilde w^2| |\sigma(W^1x+b^1)-\sigma(\tilde W^1x+\tilde b^1)|\\
&\le &  |b^2-\tilde b^2| + |w^2-\tilde w^2| \sigma(W^1x+b^1) + |\tilde w^2| |\sigma'(\cdot)| \Big(|b^1-\tilde b^1| +\|(W^1-\tilde W^1)x\| \Big) \\
&\le & \|\theta-\tilde\theta\| \max\{1,\sigma_{max}, B^\theta \sigma'_{max}, B^\theta B^x \sigma'_{max}\}.
\end{eqnarray*}
On the other hand we have
\begin{eqnarray*}
|f(\theta,x)-f(\theta,\tilde x)| &\le&  |w^2|  \sigma'_{max}  \|W^1(x-\tilde x)\| \le (B^\theta)^2 \sigma'_{max} \|x-\tilde x\|.
\end{eqnarray*}
This implies part a).

b) The gradient involves products of parameters of the neural net and $\sigma'$. More precisely we obtain:
\begin{eqnarray*}
    f_{b^1}(\theta,x) &=& w^2 \sigma'(W^1x+b^1) \\
    f_{b^2}(\theta,x) &=& 1 \\
    \nabla_{W^1} f(\theta,x) &=& w^2 \sigma'(W^1x+b^1) \cdot x \\
    f_{w^2}(\theta,x) &=&  \sigma(W^1x+b^1)
\end{eqnarray*}
and 
$$\nabla_\theta f=\Big(f_{b^1}, f_{b^2}, \nabla_{W^1} f, f_{w^2}  \Big). $$
In particular if $f_{b^1}, f_{b^2}, \nabla_{W^1} f, f_{w^2}$  are Lipschitz-continuous with constants $L_{b^2}=0$ and  $L_{b^1},L_{W^1},L_{w^2}$, then $\nabla_\theta f$ is  Lipschitz-continuous with constant  $L_f^\nabla =  L_{b^1}+ L_{b^2}+L_{W^1}+L_{w^2}.$ The Lipschitz constants can be obtained similar to part a).
\end{proof}

\begin{lemma}\label{lem:LR}
Under Assumption \ref{as:1}-\ref{ass:rFLip}
the mapping $\theta \mapsto R_T^\theta$ is a.s. locally Lipschitz continuous. On the set $\Theta\times G$ the function is globally Lipschitz continuous with a constant $K_{R,T}$ which can be obtained recursively, i.e.
$$| R_T^\theta(\omega) - R_T^{\tilde \theta}(\omega)| \le K_{R,T} \|\theta-\tilde\theta\| \mbox{ for all } \omega \in G \mbox{ and } \theta,\tilde\theta\in\Theta.$$
\end{lemma}

\begin{proof}    
We have
\begin{eqnarray*}
| R_T^\theta - R_T^{\tilde\theta}| \le \sum_{k=0}^{T-1} | r(s_k^\theta,a_k^\theta) - r(s_k^{\tilde\theta},a_k^{\tilde\theta})| \le \sum_{k=0}^{T-1} L_{r} \big( | s_k^\theta- s_k^{\tilde\theta}| +|a_k^\theta -a_k^{\tilde\theta}| \big)
\end{eqnarray*}
In what follows denote
\begin{eqnarray*}
\Delta s_t &:=& s_t^\theta- s_t^{\tilde\theta},\\
\Delta a_t &:=& a_t^\theta- a_t^{\tilde\theta},\\
\Delta R_t &:=& R_t^\theta- R_t^{\tilde\theta}.
\end{eqnarray*}
We prove now that there exists constants $K_{s,t}, K_{a,t}, K_{R,t}$ such that for $t=0,1,\ldots ,T-1$
\begin{equation}\label{eq:deltasaR} | \Delta s_t| \le K_{s,t} \|\theta-\tilde\theta\|,\quad | \Delta a_t| \le K_{a,t}\|\theta-\tilde\theta\|, \quad | \Delta R_t| \le K_{R,t}\|\theta-\tilde\theta\|
\end{equation}
which implies the statement.
We do this by induction on $t$. For $t=0$ we have with Lemma \ref{lem:NNLipschitz} that
$$ \Delta s_0=0, \quad \Delta R_0 =0, \quad \Delta a_0\le L_{f} \|\theta-\tilde\theta\|. $$
Now suppose the statement is true up to time $t-1$.
From the state dynamics, the parametrization of the action and the definition of $R_t^\theta$ we obtain
\begin{itemize}
\item[(1)] $| \Delta s_t| \le \tilde L_{F} \big( |\Delta s_{t-1}|+  |\Delta a_{t-1}|\big)$
\item[(2)] $| \Delta R_t| \le \sum_{k=0}^{t-1}  L_{r} \big( | \Delta s_k| + | \Delta a_k| \big)$.
\item[(3)] $| \Delta a_t| \le L_{f} \big( \|\theta-\tilde\theta\| + |\Delta s_{t}|+ |\Delta R_{t}|\big)$
\end{itemize}
These relations obviously imply that the statement in \eqref{eq:deltasaR} is also true for time point $t$. It is also possible to derive the following recursive equations for the Lipschitz constants for $t=1,\ldots, T$:
\begin{eqnarray*}
    K_{s,t} = \tilde L_F \big( K_{s,t-1} +K_{a,t-1}\big),\\
    K_{R,t} = L_r \sum_{k=0}^{t-1}\big( K_{s,k} +K_{a,k}\big),\\
    K_{a,t} = L_f  ( 1+ K_{s,t} +K_{R,t}).
\end{eqnarray*}
The recursion may be simplified when we use larger constants. 
%For example  by setting $K_t = \max\{K_{s,t},K_{a,t}\}.$
\end{proof}

\begin{lemma}\label{lem:LRgrad}
Under Assumption \ref{as:1}-\ref{ass:rFLipgrad}
the mapping $\theta \mapsto \nabla_\theta R_T^\theta$ is a.s. locally Lipschitz continuous.  On the set $\Theta\times G$ the function is globally Lipschitz continuous with a constant $M_{R,T}$ which can be obtained recursively, i.e.
$$| \nabla_\theta R_T^\theta(\omega) - \nabla_\theta R_T^{\tilde\theta}(\omega)| \le M_{R,T} \|\theta-\tilde\theta\| \mbox{ for all } \omega \in G \mbox{ and } \theta,\tilde\theta\in\Theta.$$
\end{lemma}

\begin{proof}
%{\em Proof of Lemma \ref{lem:LRgrad}:}
We have
\begin{eqnarray}\label{eq:nablaR}
\nabla_\theta R_T^\theta = \sum_{t=0}^{T-1}  \partial_sr(s_t^\theta,a_t^\theta) \nabla_\theta s_t^\theta + \partial_ar(s_t^{\theta},a_t^{\theta}) \nabla_\theta a_t^\theta
\end{eqnarray}
In view of Lemma \ref{lem:lipschitzproperties} we have to show that $\partial_sr, \partial_ar,s_t^\theta, a_t^\theta, \nabla_\theta s_t^\theta , \nabla_\theta a_t^\theta$ are all locally Lipschitz.  The Lipschitz properties of $\partial_sr, \partial_ar$ follow from  Assumption \ref{ass:rFLipgrad}. The Lipschitz properties of $s_t^\theta, a_t^\theta$ follow from the previous proof. It remains to show the Lipschitz property of $\nabla_\theta s_t^\theta , \nabla_\theta a_t^\theta.$ 
In what follows denote
\begin{eqnarray*}
D_{s,t}  &:=& \nabla_\theta s_{t}^\theta- \nabla_\theta s_{t}^{\tilde\theta},\\
D_{a,t} &:=& \nabla_\theta a_t^\theta- \nabla_\theta a_t^{\tilde\theta},\\
D_{R,t}  &:=& \nabla_\theta R_t^\theta- \nabla_\theta R_t^{\tilde\theta}.
\end{eqnarray*}
We prove now that there exists constants $M_{s,t}, M_{a,t}, M_{R,t}$ such that for $t=0,1,\ldots ,T-1$
\begin{equation}\label{eq:deltanablasaR} \| D_{s,t}\| \le M_{s,t} \|\theta-\tilde\theta\|,\quad \| D_{a,t}\| \le M_{a,t}\|\theta-\tilde\theta\|, \quad \| D_{R,t} \| \le M_{R,t}\|\theta-\tilde\theta\|.
\end{equation}
As in the previous proof we proceed by induction.
For $t=0$ we have with Lemma \ref{lem:NNLipschitz} that
$$ D_{s,0} =0, \quad D_{R,0} =0, \quad D_{a,0}\le L_{f}^\nabla \|\theta-\tilde\theta\|. $$
Now suppose the statement is true up to time $t-1$. 
Before we proceed with the induction step let us recall that $x=(s,R,z,t)$ and $\partial_sf := \frac{\partial}{\partial s}f$ and $\partial_Rf := \frac{\partial}{\partial R}f.$   Recall the bounds that we introduced on 
 $\Theta\times G$ in subsection \ref{sec:bounds}.

We further have (cp. subsection \ref{sec:bounds}):
\begin{eqnarray}\label{eq:nablas}
\nabla_\theta s_t^\theta &=& \partial_sF(s_{t-1}^\theta,a_{t-1}^\theta,\varepsilon_t) \nabla_\theta s_{t-1}^\theta + \partial_aF(s_{t-1}^\theta,a_{t-1}^\theta,\varepsilon_t) \nabla_\theta a_{t-1}^\theta \\\label{eq:nablaa2}
\nabla_\theta a_t^\theta &=& f_\theta(\theta,s_{t}^\theta,R_{t}^\theta,z_t,t)+ \partial_s f(\theta,s_{t}^\theta,R_{t}^\theta,z_t,t) \nabla_\theta s_{t}^\theta + \partial_R f(\theta,s_{t}^\theta,R_{t}^\theta,z_t,t) \nabla_\theta R_{t}^\theta.
\end{eqnarray}
From the first equation \eqref{eq:nablas} we obtain 
\begin{eqnarray*}
    M_{s,t} &=& B_F^s M_{s,t-1} + B_\nabla^s L_\nabla^{F,s}(K_{s,t-1} +K_{a,t-1})\\
    && + B_F^a M_{a,t-1} + B_\nabla^a L_\nabla^{F,a}(K_{s,t-1} +K_{a,t-1}).
\end{eqnarray*}

From the equation for $\nabla_\theta R_t^\theta$ in \eqref{eq:nablaR} we obtain:

$$M_{R,t} = \sum_{k=0}^{t-1} B_r^s M_{s,k} + B_r^a M_{a,k} + \big( B_\nabla^s L_\nabla^{r,s}+B_\nabla^a L_\nabla^{r,s} \big)\big(K_{s,k}+K_{a,k} \big)$$
It remains to discuss $M_{a,t}.$ In order to do this we need a further look at $\partial_if, i=s,a$ compare \eqref{eq:definition_NN_formula}:
\begin{eqnarray*}
    \partial_sf(\theta,x) &=& w^2 \sigma(w^1_s s + w^1_R R+ w^1_z z + w^1_t (T-t)+ b^1) w^1_s\\ 
    \partial_Rf(\theta,x) &=& w^2 \sigma(w^1_s s + w^1_R R+ w^1_z z + w^1_t (T-t)+ b^1) w^1_R.
\end{eqnarray*}
Using these formulas we obtain for the first expression (the second expression is similar):
\begin{eqnarray*}
  &&  |\partial_sf(\theta,s_t^\theta,R_t^\theta,z_t,t)-\partial_sf(\tilde\theta,s_t^{\tilde\theta},R_t^{\tilde\theta},z_t,t)| \\
&\le& |\partial_sf(\theta,s_t^\theta,R_t^\theta,z_t,t)-\partial_sf(\tilde\theta,s_t^{\theta},R_t^{\theta},z_t,t)| + |\partial_sf(\tilde\theta,s_t^\theta,R_t^\theta,z_t,t)-\partial_sf(\tilde\theta,s_t^{\tilde\theta},R_t^{\tilde\theta},z_t,t)|
\end{eqnarray*}
For the first expression we obtain:
\begin{eqnarray*}
  &&  |\partial_sf(\theta,s_t^\theta,R_t^\theta,z_t,T)-\partial_sf(\tilde\theta,s_t^{\theta},R_t^{\theta},z_t,T)| \le  \|\theta-\tilde\theta\|\;  \max\{B^\theta \sigma'_{max}, (B^\theta)^2B^x \sigma''_{max}  \}
\end{eqnarray*}
For the second expression we obtain:
\begin{eqnarray*}
  &&  |\partial_sf(\tilde\theta,s_t^\theta,R_t^\theta,z_t,T)-\partial_sf(\tilde\theta,s_t^{\tilde\theta},R_t^{\tilde\theta},z_t,T)| \le  \|\theta-\tilde\theta\|\;   (B^\theta)^3 \sigma''_{max} (K_{s,t}+K_{R,t})
%  &\le &  \Big\{2 B^\theta \sigma'_{max} +(B^\theta)^2 \sigma''_{max}  \Big(B^s +B^\theta K_{s,t} +B^R +B^\theta K_{R,T} +B^z +T \Big) \Big\} =: \|\theta-\tilde\theta\| L_f^{\nabla_x}
\end{eqnarray*}
In total this yields 
\begin{eqnarray*}
  &&  |\partial_sf(\theta,s_t^\theta,R_t^\theta,z_t,T)-\partial_sf(\tilde\theta,s_t^{\tilde\theta},R_t^{\tilde\theta},z_t,T)| \\
&\le& \|\theta-\tilde\theta\|\; \max\{(B^\theta)^3 \sigma''_{max} (K_{s,t}+K_{R,t}), B^\theta \sigma'_{max}, (B^\theta)^2B^x \sigma''_{max} \  \}\\
&=:& \|\theta-\tilde\theta\|\;  L_{f,t} ^{\nabla_x} 
\end{eqnarray*}

The Lipschitz constants for $\partial_Rf$ is the same.
Thus, we obtain from \eqref{eq:nablaa2} that
$$ M_{a,t} = L_f^\nabla+ B_\nabla^{f,s} M_{s,t} + B_\nabla^s L_{f,t} ^{\nabla_x} +
 B_\nabla^{f,R} M_{R,t} + B_\nabla^R L_{f,t} ^{\nabla_x} $$
Finally, the induction step is complete and the statement follows.
\end{proof}

\begin{proof} of Theorem \ref{lem:Lipschitz_gradient}:
First note that we can write the gradient more explicitly as
\begin{eqnarray*}
     \nabla_\theta \mathcal{L}_L(\theta) &=&
   2 \sum_{\ell=1}^L \beta_\ell  \left[ \left(Re(\varphi^\star)(u_\ell)-\EE_G[\cos(u_\ell R_T^\theta)]\right)\cdot \mathbb{E}_G\left[ \sin(u_\ell R_T^\theta) \nabla_\theta R_T^\theta \right] \right] \\
   && - 2 \sum_{\ell=1}^L  \beta_\ell \left[ \left(Im(\varphi^\star)(u_\ell)-\EE_G[\sin(u_\ell R_T^\theta)]\right)\cdot \mathbb{E}_G\left[ \cos(u_\ell R_T^\theta) \nabla_\theta R_T^\theta \right] \right] 
\end{eqnarray*}
Now from Lemma \ref{lem:LR} it follows that
$$\theta \mapsto \EE_G[\cos(u R_T^\theta)] \quad\mbox{   and    } \quad\theta \mapsto \EE_G[\sin(u R_T^\theta)]$$
are both Lipschitz-continuous with constant $u K_{R,T}.$
Lemma \ref{lem:LR} together with Lemma \ref{lem:LRgrad} imply that 
$$\theta \mapsto \EE_G[\sin(u R_T^\theta) \nabla_\theta R_T^\theta ] \quad \mbox{   and    }\quad \theta \mapsto \EE_G[\cos(u R_T^\theta) \nabla_\theta R_T^\theta]$$
are both Lipschitz-continuous with constant $M_{R,T} + u K_{R,T} B_\nabla^R.$ Using the rules for the Lipschitz constants from Lemma \ref{lem:lipschitzproperties} we obtain 
the following Lipschitz constant for $ \nabla_\theta \mathcal{L}_L$ using the $L_1$-norm:
\begin{eqnarray*}
    \tilde L &=& 8 M_{R,T} \sum_{\ell=1}^L \beta_\ell  + 12 B_\nabla^R K_{R,T} \sum_{\ell=1}^L |u_\ell\beta_\ell|.
\end{eqnarray*}
Noting that for $x\in\RR^n$ we have $\|x\|_2 \le \|x\|_1\le \sqrt{n} \|x\|_2$ we obtain 
 the result with $\mathbf{L} =\sqrt{7} \tilde L$.
\end{proof}

\begin{proof}
of Lemma \ref{lem:bias}: We fix an iteration index $k\in\mathbb N$ and work throughout \emph{conditionally} on the
$\sigma$--algebra $\cF_k^G$.
Under this conditioning, the parameter $\theta_k$ is deterministic, and all bounds derived
on the good event $G$ hold pathwise.
In particular, by Lemma \ref{lem:bound1}
\[
\|\nabla_\theta R_T^{\theta_k}\|\;\le\; B_\nabla^R
\qquad\text{almost surely on } G.
\]

%For each frequency node $\omega_\ell$, 
Let $(R_m,\nabla_\theta R_m)_{m=1}^M$ be i.i.d.\ copies of
$(R_T^{\theta_k},\nabla_\theta R_T^{\theta_k})$ under the current parameter $\theta_k$. Thus, these Monte Carlo samples are conditionally i.i.d.\ given $\cF_k^G$. All expectations below are therefore taken with respect to these samples only, while $\theta_k$ is treated as fixed.

\medskip
\noindent\textbf{Step 1: Bias.}

Fix a frequency node $u=u_\ell$ and, to simplify notation, suppress the index $\ell$
throughout this step.
Recall that for this node the empirical characteristic function is given by
\[
\widehat\varphi(u)
:=
\frac1M\sum_{j=1}^M e^{iu R_j}.
\]

For each sample index $j\in\{1,\dots,M\}$, define the pathwise gradient contribution
\[
U_j(u)
:=
iu\,e^{iu R_j}\,\nabla_\theta R_j
\;\in\;\mathbb{C}^7.
\]
With this notation, the contribution of the frequency node $u_\ell$ to the full estimator
$g(\theta_k)$ can be written as
\[
g_\ell
=
\frac{2\beta_\ell}{M}\sum_{j=1}^M
\text{Re}\Big(
\overline{\widehat\varphi(u)-\varphi^\star(u)}\;
U_j(u)
\Big).
\]

We emphasize that the same Monte Carlo batch
$(R_m)_{m=1}^M$ is used both to construct the empirical characteristic function
$\widehat\varphi(u)$ and to evaluate the gradient terms $U_j(u)$.
This coupling is the source of the finite-sample bias analyzed below. Taking conditional expectation and using linearity of $\text{Re}(\cdot)$,
\[
\EE[g_\ell\mid \cF_k^G]
=
\frac{2\beta_\ell}{M}\sum_{j=1}^M
\text{Re}\Big(\EE[\overline{\widehat\varphi(u)}\,U_j(u)\mid \cF_k^G]
-\overline{\varphi^\star(u)}\,\EE[U_j(u)\mid \cF_k^G]\Big).
\]
Because the samples are i.i.d., it suffices to compute $\EE[\overline{\widehat\varphi(u)}\,U_1(u)\mid \cF_k^G]$:
\[
\EE[\overline{\widehat\varphi(u)}\,U_1(u)\mid \cF_k^G]
=
\frac1M\EE[\overline{e^{iu R_1}}\,U_1(u)\mid \cF_k^G]
+\frac{M-1}{M}\EE[\overline{e^{iu R_2}}\mid \cF_k^G]\EE[U_1(u)\mid \cF_k^G].
\]
Now $\overline{e^{iu R_1}}\,U_1(u)= iu \nabla_\theta R_1$,
and we denote $\EE[{e^{iu R_2}}\mid \cF_k^G]=\EE_G[{e^{iu R_T^{\theta_k}}}] =:{\varphi_{\theta_k}(u)}$.
Thus $\EE[U_1(u)\mid \cF_k^G]=\nabla_\theta\varphi_{\theta_k}(u)$
and we obtain
\[
\EE[\overline{\widehat\varphi(u)}\,U_1(u)\mid \cF_k^G]
=
\frac1M\,iu\,\EE[\nabla_\theta R_1\mid \cF_k^G]
+\frac{M-1}{M}\,\overline{\varphi_{\theta_k}(u)}\,\nabla_\theta\varphi_{\theta_k}(u).
\]
Plugging back yields
\[
\EE[g_\ell\mid \cF_k^G]
=
2\beta_\ell\,\text{Re}\Big(\overline{\varphi_{\theta_k}(u)-\varphi^\star(u)}\,\nabla_\theta\varphi_{\theta_k}(u)\Big)
+\frac{2\beta_\ell}{M}\text{Re}\Big(iu\,\EE[\nabla_\theta R_1\mid \cF_k^G]-\overline{\varphi_{\theta_k}(u)}\,\nabla_\theta\varphi_{\theta_k}(u)\Big).
\]
Summing over $\ell$ gives \eqref{eq:bias-decomposition}--\eqref{eq:bias-term}.

To bound the bias, use $\|\nabla_\theta R_1\|_2\le \|\nabla_\theta R_1\|\le B_\nabla^R$ and
$\|\nabla_\theta\varphi_{\theta_k}(u)\|_2
\le \|\nabla_\theta\varphi_{\theta_k}(u)\|
\le \EE[|u|\,\|\nabla_\theta R_1\|\mid \cF_k^G ]\le |u| B_\nabla^R$,
and $|\varphi_{\theta_k}(u)|\le 1$, obtaining \eqref{eq:bias-bound}.

\medskip
\noindent\textbf{Step 2: Second moment.}
For each $(\ell,j)$ define
\[
Y_{\ell,j}
:=\mathrm{Re}\Big(\overline{\hat{\varphi}(u_\ell)-\varphi^*(u_\ell)}\,iu_\ell e^{iu_\ell R_j}\nabla_\theta R_j\Big)\in\RR^7,
\qquad
g(\theta_k)=\frac{2}{M}\sum_{\ell=1}^{L}\beta_\ell\sum_{j=1}^M Y_{\ell,j}.
\]
We have
$|\hat{\varphi}(u_\ell)-\varphi^*(u_\ell)|\le  2$
and $|e^{i\cdot}|=1$. Hence, for all $(\ell,j)$ we obtain on  $G$
\begin{equation}\label{eq:Y-bound}
\|Y_{\ell,j}\|
\le 2|u_\ell|\,\|\nabla_\theta R_j\|
\le 2|u_\ell|\,B_\nabla^R.
\end{equation}

Next, by the triangle inequality,
\[
\Big\|\sum_{j=1}^M Y_{\ell,j}\Big\|
\le \sum_{j=1}^M \|Y_{\ell,j}\|
\le M \max_{1\le j\le M}\|Y_{\ell,j}\|.
\]
%Squaring and using \eqref{eq:Y-bound} yields the deterministic estimate, valid on $\cG_\delta^{(T)}$,
%\begin{equation}\label{eq:sumj-bound}
%\Big\|\sum_{j=1}^M Y_{\ell,j}\Big\|^2
%\le M^2 \big(2|u_\ell|K_{\mathrm{grad}}^{(\delta)}\big)^2
%= 4M^2u_\ell^2\big(K_{\mathrm{grad}}^{(\delta)}\big)^2.
%\end{equation}
%Plugging \eqref{eq:sumj-bound} into the definition of $g(\theta_k)$ and using again the triangle inequality,
Thus, we obatin
\[
\|g(\theta_k)\|
\le \frac{2}{M}\sum_{\ell=1}^{L}|\beta_\ell|
\Big\|\sum_{j=1}^M Y_{\ell,j}\Big\|
\le \frac{2}{M}\sum_{\ell=1}^{L}|\beta_\ell|\,
\Big(2M|u_\ell|B_\nabla^R\Big)
=4B_\nabla^R\sum_{\ell=1}^{L}|\beta_\ell u_\ell|.
\]
Therefore, on $G$,
\begin{equation}\label{eq:g-det-bound}
\|g(\theta_k)\|_2^2
\le 16\big(B_\nabla^R\big)^2
\Big(\sum_{\ell=1}^{L}|\beta_\ell u_\ell|\Big)^2.
\end{equation}
Taking conditional expectations preserves the inequality, so we obtain the statement.
\end{proof}

\bibliographystyle{apalike}
\bibliography{sample}
 
\end{document}
